% 本文件由 LaTeX 排版 Agent 根据有效 translation.json 自动生成。
% author=Hermas; work_id=0201; title=《牧者》

\chapter{《牧者》}

% structure: p0001 source=h1 target=omitted reason=page-title-already-rendered-by-parent

% structure: p0002 source=h2 target=section reason=relative-heading-rank; base=h2

\section{第一书：神视}

\begin{itemize}
\item 第一神视——反对污秽与骄傲的思念，以及赫尔马斯在管教其子上的疏忽。
\end{itemize}

\begin{itemize}
\item 第二神视——再次论及他管教其多言妻子和贪欲儿子上的疏忽，以及他的品性。
\end{itemize}

\begin{itemize}
\item 第三神视——论凯旋教会的建造，以及各种被弃绝之人的类别。
\end{itemize}

\begin{itemize}
\item 第四神视——论将要临到人身上的试炼与磨难。
\end{itemize}

\begin{itemize}
\item 第五神视——论诫命。
\end{itemize}

% structure: p0008 source=h2 target=section reason=relative-heading-rank; base=h2

\section{第二书：诫命}

\begin{itemize}
\item 第一条诫命——论对天主的信德。
\end{itemize}

\begin{itemize}
\item 第二条诫命——论避免恶言，以及以纯朴施舍。
\end{itemize}

\begin{itemize}
\item 第三条诫命——论避免谎言，以及赫尔马斯为其伪善所作的忏悔。
\end{itemize}

\begin{itemize}
\item 第四条诫命——论因奸淫休妻。
\end{itemize}

\begin{itemize}
\item 第五条诫命——论心中的忧愁与忍耐。
\end{itemize}

\begin{itemize}
\item 第六条诫命——如何认识每个人身旁的两个神灵，以及如何区分二者的提示。
\end{itemize}

\begin{itemize}
\item 第七条诫命——论敬畏天主，而不畏惧魔鬼。
\end{itemize}

\begin{itemize}
\item 第八条诫命——我们应当避开邪恶，而行善。
\end{itemize}

\begin{itemize}
\item 第九条诫命——必须不停息、以坚定不移的信心向天主祈祷。
\end{itemize}

\begin{itemize}
\item 第十条诫命——论忧愁，以及不可使在我们内的圣神忧愁。
\end{itemize}

\begin{itemize}
\item 第十一条诫命——圣神与先知须凭其行为考验；亦论及两种神类。
\end{itemize}

\begin{itemize}
\item 第十二条诫命——论双重欲望。天主的诫命可以遵守，信徒不应畏惧魔鬼。
\end{itemize}

% structure: p0021 source=h2 target=section reason=relative-heading-rank; base=h2

\section{第三书：比喻}

\begin{itemize}
\item 第一个比喻——因在此世没有永存的城邑，我们当寻求那将来的。
\end{itemize}

\begin{itemize}
\item 第二个比喻——如同葡萄树依靠榆树，富人也因穷人的祈祷得到扶助。
\end{itemize}

\begin{itemize}
\item 第三个比喻——如同冬天绿树与枯树无法分辨，同样在此世义人与不义之人也无法分辨。
\end{itemize}

\begin{itemize}
\item 第四个比喻——如同夏天有生命的树凭果实与绿叶与枯树区分，同样在来世义人与不义之人将在福乐中分别。
\end{itemize}

\begin{itemize}
\item 第五个比喻——论真正的斋戒及其赏报；亦论身体的纯洁。
\end{itemize}

\begin{itemize}
\item 第六个比喻——论两类纵欲之人，及其死亡、堕落与惩罚的期限。
\end{itemize}

\begin{itemize}
\item 第七个比喻——忏悔的人必须结出与忏悔相称的果实。
\end{itemize}

\begin{itemize}
\item 第八个比喻——被选者与忏悔者的罪恶有多种，但所有人将按其忏悔与善行的程度得到赏报。
\end{itemize}

\begin{itemize}
\item 第九个比喻——在争战与凯旋教会的建造中的伟大奥秘，
\end{itemize}

\begin{itemize}
\item 第十个比喻——论忏悔与施舍。
\end{itemize}

\SourceRecord{Translated by F. Crombie.；Ante-Nicene Fathers；Vol. 2.；Edited by Alexander Roberts, James Donaldson, and A. Cleveland Coxe.；Buffalo, NY: Christian Literature Publishing Co.,；1885.；<http://www.newadvent.org/fathers/0201.htm>.}

% structure: p0001 source=h1 target=section reason=page-title

\section{\WorkTitle{赫耳玛斯牧者（卷一）}}

% structure: p0002 source=h2 target=subsection reason=relative-heading-rank; base=h2

\subsection{神视一}

\Emph{反对污秽和骄傲的思想，以及赫耳玛斯在管教其子上的疏忽。}\par

% structure: p0004 source=h3 target=subsubsection reason=relative-heading-rank; base=h2

\subsubsection{第一章}

那养育我的人，把我卖给罗马一个名叫罗德的人。多年之后，我认出了她，并开始像爱姐妹一样爱她。过了一些时候，我看见她在台伯河中沐浴；我便伸手拉她上岸。她的美貌使我想道：\Quote{我若能娶到像她那样美丽善良的妻子，该是多么幸福的人啊。}这念头只是在我心中闪过，仅此而已。不久之后，我走在通往乡村的路上，赞美天主的受造物，思想它们是多么宏伟、美丽、有力，我便睡着了。圣神将我提去，带到一个无路之处，人不能行走，因为它位于岩石之中；因水势而崎岖难行。过了这条河，我来到一片平原。于是屈膝跪下，开始向主祈祷，并忏悔我的罪。我祈祷时，天开了，我看见我所渴慕的那位女子从天上向我致意，说：\Quote{万福，赫耳玛斯！}我抬头望着她，说：\Quote{夫人，您在这里做什么？}她回答我说：\Quote{我被带到此处，要在主前控告你的罪。}\Quote{夫人，}我说，\Quote{难道您要成为控告我的对象吗？}\Quote{不，}她说，\Quote{但请听我要对你说的话。那居住在天上、从无中造出万有、并为其圣教会使之繁衍生长的天主，因你得罪了我而向你发怒。}我回答她：\Quote{夫人，我得罪了您吗？如何得罪的？何时说了不当的话？我不是一向把您看作夫人吗？不是一向像尊重姐妹一样尊重您吗？您为何诬告我犯此邪恶和不洁？}她微笑着回答我说：\Quote{邪恶的欲望已在你的心中升起。你难道不认为，当义人心中生出邪念时，他便是犯罪吗？在这种情况下便有罪，而且罪很大，}她说，\Quote{因为义人的思念本应是正义的。藉着正义的思念，他的品性在诸天中得以建立，主在一切事务上都对他仁慈。但那些心怀恶念的人，是给自己招致死亡和束缚；尤其是那些贪恋现世、以财富为荣、不展望来生福乐的人。他们的懊悔必将众多，因为他们没有希望，已对自己和生命绝望。但你要向天主祈祷，祂必医治你的罪，以及你全家和所有圣人的罪。}\par

% structure: p0006 source=h3 target=subsubsection reason=relative-heading-rank; base=h2

\subsubsection{第二章}

她说完这些话，天就关了。我充满忧愁和恐惧，对自己说：\Quote{若这罪归在我身上，我如何能得救？或如何能为我这极重的罪向天主赎罪？我该用什么话来求主怜悯我？}我正思想并琢磨这些事时，看见对面有一把白色的大椅子，是用白羊毛做的。接着来了一位老妇人，穿着华美的长袍，手里拿着一本书；她独自坐下，向我致意说：\Quote{万福，赫耳玛斯！}我悲伤地流着泪对她说：\Quote{夫人，万福！}她对我说：\Quote{赫耳玛斯，你为何如此沮丧？你一向忍耐、节制，总是面带微笑。为何这样愁眉不展，不快活？}我回答她说：\Quote{啊，夫人，我被一位极好的妇人责备，她说我得罪了她。}她说：\Quote{天主的仆人绝不可行此事。但或许你心中对她起了欲念。在事奉天主的人身上，这样的念头便产生罪。因为在极其贞洁且受过考验的灵里，欲念一件恶事是邪恶可怕的；尤其赫耳玛斯，他素来远离一切邪念，充满纯朴和极大的天真无邪。}\par

% structure: p0008 source=h3 target=subsubsection reason=relative-heading-rank; base=h2

\subsubsection{第三章}

\Quote{但天主并非因此对你发怒，而是要你归化你的家，他们曾得罪主并得罪你，他们的父母。虽然你爱你的儿子们，但你并没有警告你的家，反而容许他们严重败坏。因此主向你发怒，但祂必医治你家中所有邪恶。因为由于他们的罪和不义，你因这世界的事务而毁败。但现在主的仁慈已经怜悯你和你的家，并将坚固你，将你安立在祂的荣耀中。只是你不要轻易动摇，而要鼓起勇气，安抚你的家。因为正如铁匠锤出他的作品，成就他所愿的一切，同样，正义的日常言语将胜过一切不义。因此不要停止劝诫你的儿子们；因为我知道，如果他们全心悔改，他们将与圣人一同被记在生命册上。}说完这些话，她对我说：\Quote{你愿意听我宣读吗？}我对她说：\Quote{夫人，我愿意。}\Quote{那么请听，侧耳听天主的荣耀。}然后我听见了从她而来的话，宏伟而奇妙，我的记忆不能保留。因为所有的话都可怕，是人不能承受的。然而最后的话我记得，因为它们对我们有益且温和。\Quote{看！万军的天主，用祂不可见的强大能力和大智慧创造了世界，用祂荣耀的谋划以美丽环绕祂的创造，用祂强有力的言语安定诸天，在水上立定大地根基，并用自己的智慧和眷顾创造了祂神圣的教会，祂所祝福的。看！祂挪移诸天和山岭、丘陵和海洋，万象在祂的选民面前变得平坦，为将祂所应许的祝福赐给他们，带着极大的荣耀和喜乐，只要他们持守所领受的天主诫命，怀着大信心。}\par

% structure: p0010 source=h3 target=subsubsection reason=relative-heading-rank; base=h2

\subsubsection{第四章}

她读完后，从椅子上站起来，四个年轻人前来抬走椅子，往东边去了。她叫我到她面前，触摸我的胸膛，对我说：\Quote{你喜欢我的宣读吗？}我对她说：\Quote{夫人，末尾的话使我喜悦，但开头的话残酷严苛。}然后她对我说：\Quote{末尾的话是为义人；开头的话是为外邦人和叛教者。}她对我说话时，出现两个人，把她扛在肩上，往椅子所在的东西方向去了。她带着喜悦的面容离去了；她边走边对我说：\Quote{要像个男子汉，\Person{赫尔玛}。}\par

% structure: p0012 source=h2 target=subsection reason=relative-heading-rank; base=h2

\subsection{第二神视}

\Emph{再次，论他在惩罚其多言妻子及好色儿子上的疏忽，及其品德。}\par

% structure: p0014 source=h3 target=subsubsection reason=relative-heading-rank; base=h2

\subsubsection{第一章}

大约与前一年同一时间，我前往乡下，途中回忆起当年的神视。圣神再次提携我，带我到去年曾去过的同一地方。到达那地方，我屈膝跪地，开始向主祈祷，赞美祂的名，因祂认为我配得，并向我启示了我从前的罪过。祈祷后站起来，我看见对面有位老妇人，就是去年我所见到的那位，正走着读着一本书。她对我说：\Quote{你能将这些事报告给天主的选民吗？}我对她说：\Quote{夫人，我无法尽记于心，但请把书给我，我将誊抄下来。}她说：\Quote{拿去吧，以后要还给我。}于是我接过书，来到乡间一处，从头到尾逐字誊抄；但我未能明白其音节。然而我一抄完那书，它突然从我手中被夺去；但我没有看见夺书的人。\par

% structure: p0016 source=h3 target=subsubsection reason=relative-heading-rank; base=h2

\subsubsection{第二章}

十五天后，在我禁食并向主多多祈祷后，那些文辞的意义向我启示了。这文辞的内容如下：\Quote{你的后裔，\Person{赫尔玛}，得罪了天主，并亵渎了主，且因他们的大恶背叛了父母。他们身为父母的叛徒，却未从背叛中获利。如今他们又加添了情欲和不洁的玷污，于是他们的不义便满了。你要将这些话宣告给你的众子女和你的妻子——她将成为你的姊妹。因为她没有勒住舌头，以之言行不义；但听了这些话，她会自制，并得到怜悯。因为在你将主命我启示给你的这些话宣告给他们之后，他们从前所犯的一切罪过都将被赦免，并且所有至今仍犯罪的圣人，若他们全心悔改，驱除心中一切疑虑，也将获得宽恕。因为主已指着祂的荣耀为祂的选民起誓：若他们在某一指定之日以后犯罪，必不得救。因为义人的悔改有其限度。所有圣人的悔改的日子已经满了；但对外邦人，悔改可能直到末日。因此，你要告诉那些管理教会的人，叫他们以正义引导自己的道路，使他们得以带着极大荣耀完全领受所应许的。所以，行义的人哪，要站得稳固，不要疑惑，好使你们的旅程与圣天使同在。有福的是你们，忍受那将要来临的大苦难，有福的是那不肯否认自己生命的人。因为主已指着圣子起誓：那些否认他们主的人已经绝望地抛弃了自己的生命，因为甚至现在，在将来临的日子里，这些人将要否认祂。对于早期否认的人，因祂极其温柔的仁慈，天主已成为仁慈的。}\par

% structure: p0018 source=h3 target=subsubsection reason=relative-heading-rank; base=h2

\subsubsection{第三章}

\Quote{但是，\Person{赫尔玛}，你不要记念你儿女们对你的不义，也不要忽略你的姊妹，好叫他们从先前的罪中得洁净。因为他们若蒙受正义的教训，你若不忘怀他们对你的不义，他们就会得洁净。因为不忘不义培植死亡。而你呢，\Person{赫尔玛}，因着你家的过犯，你承受了极大的个人磨难，因为你没有留意他们，反而粗心大意、从事邪恶的交易。但你得救了，因为你没有离弃生活的天主，也因你的纯朴和大节制。这些救了你，只要你坚定不移。它们也要拯救所有这样行事、在无欺诈和纯朴中行走的人。拥有这些美德的人必胜过一切邪恶的形态，并要存留直到永生。凡行正义的人是有福的，因为他们永不灭亡。现在你要告诉\Person{马克西穆斯}：看啊，磨难临到了。若你以为好，就再次否认。主临近那些归向祂的人，正如在\WorkTitle{埃尔达和莫达特}中所记载的，他们在旷野中向百姓说预言。}\par

% structure: p0020 source=h3 target=subsubsection reason=relative-heading-rank; base=h2

\subsubsection{第四章}

如今，我的弟兄们，我在睡梦中得了一个启示，是一位容貌俊美的少年对我说的，他说：\Quote{你以为给你那本书的老妇人是谁？}我说：\Quote{是\Person{西彼尔}。}\Quote{你错了，}他说，\Quote{那不是\Person{西彼尔}。}\Quote{那是谁呢？}我说。他说：\Quote{那是教会。}我对他说：\Quote{那么她为什么是一位老妇人？}\Quote{因为，}他说，\Quote{她是首先被造的。因此她是年老的。世界也是为了她而造的。}此后，我在家中看见了一个神视，那位老妇人来了，问我是否已经把书交给了长老。我说还没有。于是她说：\Quote{你做得很好，因为我还有些话要加上。但当我完成所有的话时，所有蒙选的人就会经由你认识它们。因此你要写两卷书，一卷送给\Person{克莱孟}，另一卷送给\Person{格拉普特}。\Person{克莱孟}要把他的送到外国去，因为已准许他这样做。\Person{格拉普特}要劝诫寡妇和孤儿。而你，要在这城中与治理教会的长老一同宣读这些话。}\par

% structure: p0022 source=h2 target=subsection reason=relative-heading-rank; base=h2

\subsection{第三神视}

\Emph{论得胜教会的建造，以及各类被弃绝的人。}\par

% structure: p0024 source=h3 target=subsubsection reason=relative-heading-rank; base=h2

\subsubsection{第一章}

我所见的神视，我的弟兄们，如下。我多次禁食，并祈求主，使祂向我启示祂应许透过那位老妇人要赐给我的启示。当夜，那位老妇人向我显现，对我说：\Quote{你既然如此渴望且切愿知道一切，就到你居住的那地方去吧；大约在第五时辰，我将会向你显现，并指示你应当看见的一切。}我问她说：\Quote{夫人，我该到哪地方去呢？}她说：\Quote{到你愿意的任何地方。}于是我便选了一处合适的地方，退下了。然而，在我开始说话并提及那地方之前，她对我说：\Quote{我会到你愿我去的地方。}于是，我便到了乡下，数着时辰，来到了我答应与她相见的地方。我见了一张象牙椅子已预备好，上面有麻布垫子，麻布垫子上又铺着细麻布罩单。看见这些陈设，却无人在那里，我开始感到敬畏，仿佛一阵战栗抓住了我，我的毛发竖立，恐惧临到我，因我见自己孤身一人。但当我回过神来，想起天主的荣耀，我便鼓起勇气，屈膝跪下，又像先前一样向天主认我的罪。随后，那老妇人走近，由我先前看见过的六个青年陪伴；她站在我身后，听我祈祷并向主认罪。她抚着我说：\Quote{\Person{赫尔玛}，停止不断为你的罪祈祷吧；要为正义祈祷，好叫你能立即在你的家中分得一份。}于是，她拉着我的手，领我到椅子前，对青年们说：\Quote{去建造吧。}青年们走了，只剩下我们两人，她对我说：\Quote{坐在这里。}我对她说：\Quote{夫人，请让我的长老们先坐。}\Quote{照我吩咐你的做，}她说，\Quote{坐下。}当我想要坐在她右边时，她不允许，反而用手势示意我坐在左边。对此我心中思忖，因她不让我坐在右边而懊恼，她说：\Quote{\Person{赫尔玛}，你懊恼吗？右边的位置是留给那些已经取悦了天主，并为祂的名受过苦难的人；而你还有许多事要做，才能与他们同坐。但你现在要持守你的纯朴，你就会与他们同坐，并与所有行他们所行、忍受他们所忍受的人同坐。}\par

% structure: p0026 source=h3 target=subsubsection reason=relative-heading-rank; base=h2

\subsubsection{第二章}

\Quote{他们承受了什么？}我说。她回答：\Quote{听着：为天主的名，他们承受了鞭打、监禁、巨大的苦难、十字架、野兽。为此，圣化的份被分派给他们坐在右边，并给每一个将为天主的名受苦的人。至于其余的人，则被分派坐在左边。但无论是坐在右边还是左边的人，都有同样的恩赐和应许；只是坐在右边的人拥有一些光荣。那么，你渴望与他们一同坐在右边，但你的缺失很多。然而，你将从你的缺失中被洁净；所有不陷于疑惑的人，直到今日都将从他们的一切不义中被洁净。}说完这话，她想要离开。但我伏在她的脚前，因主恳求她，把她应许要给我看的异象显示给我。于是她又拉住我的手，扶起我，让我坐在左边的座位上；然后举起一根光彩夺目的杖，对我说：\Quote{你看见什么大事吗？}我说：\Quote{夫人，我什么也看不见。}她对我说：\Quote{看哪！你没看见你对面有一座建在水上的大塔，由光彩夺目的方形石头砌成吗？}因为那塔是由与她同来的那六个青年人砌成方形的。无数的人向它搬运石头，有的从深处拖来，有的从陆地运来，并交给那六个青年人。他们接过来进行建筑；那些从深处拖出来的石头，他们就地放进建筑物里：因为它们已经磨光并与其他石头精确贴合，彼此结合得如此紧密，以至于接缝都看不出来。这样，塔的建造看起来就像由一块石头做成。然而，那些从陆地取来的石头却遭遇了不同的命运：因为青年人有的拒绝了一些，有的放进了建筑，有的则砍削后远远地扔离塔。还有许多其他石头躺在塔周围，青年们没有用它们来建造；因为有些是粗糙的，有些有裂缝，有些被切得太短，有些是白色而圆的，但不适合塔的建筑。此外，我看见另一些石头被远远地扔离塔，落在公共道路上；但它们没有留在路上，而是滚到了一个无路的地方。我又看见别的石头落在火里焚烧，别的落近水边，却无法滚入水中，尽管它们想滚下去进入水中。\par

% structure: p0028 source=h3 target=subsubsection reason=relative-heading-rank; base=h2

\subsubsection{第三章}

在向我显示这些异象后，她想要离开。我对她说：\Quote{我看见这一切却不知道它是什么意思，又有什么用处呢？}她对我说：\Quote{你是个狡猾的家伙，想知道一切与塔有关的事。}\Quote{正是如此，夫人，}我说，\Quote{好让我告诉我的弟兄们，他们听了这事，就能在极大的光荣中认识主。}她说：\Quote{的确，许多人会听见，听见了，有的欢喜，有的哭泣。但即使这些人，如果他们听见并悔改，也会欢喜。那么，听这塔的比喻吧；因为我要把一切都启示给你，在启示的事上不要再给我添麻烦：因为这些启示有终结，它们已经完成了。但你还不会停止祈求启示，因为你不知羞耻。你所看见正在建造的塔，就是我自己——教会，我现在和从前向你显现。那么，关于塔，你尽管问什么，我都启示给你，好使你与圣人一同欢喜。}我对她说：\Quote{夫人，既然你这次愿意把一切都启示给我，就启示吧。}她对我说：\Quote{凡应当启示的，必将被启示；只要你的心与天主同在，无论你看见什么，都不要怀疑。}我问她：\Quote{夫人，为什么塔建在水上？}她回答：\Quote{我以前告诉过你，你还仔细追问；因此，你追问就会找到真理。那么，请听为什么塔建在水上。因为你的生命曾藉着水得救，将来也要藉着水得救。因为塔是奠基于全能而光荣的圣名之言，并由主不可见的能力维系在一起。}\par

% structure: p0030 source=h3 target=subsubsection reason=relative-heading-rank; base=h2

\subsubsection{第四章}

我回答她说：\Quote{这真是宏伟而奇妙。但那六个从事建造的青年人是谁？}她说：\Quote{他们是天主的圣天使，是最先被造的，主将祂的全部受造物交给了他们，使他们增加、建造并管理整个受造物。藉着他们，塔的建造将完成。}\Quote{但那些搬运石头的其他人是谁？}\Quote{这些也是主的圣天使，但前六个比这些更卓越。塔的建造将完成，所有的人都将围绕塔一同欢喜，并光荣天主，因为塔完成了。}我问她，说：\Quote{夫人，我想知道那些石头的遭遇如何，以及各种石头意味着什么。}她回答我说：\Quote{并不是因为你比众人更配得这启示——因为在你之前还有别的人，比你更好，这些异象本应启示给他们——而是为了天主的名字得到光荣，这启示才给了你，并且因那些心中疑惑、思忖这些事是否会发生的人，也将给予启示。告诉他们，这一切都是真的，没有一件超越真理。它们全都坚定可靠，建立在坚固的根基上。}\par

% structure: p0032 source=h3 target=subsubsection reason=relative-heading-rank; base=h2

\subsubsection{第五章}

\Quote{现在关于建筑中的石头，请听。那些纯白且彼此吻合的方形石头，是宗徒、主教、教师和执事，他们度着虔诚纯洁的生活，并贞洁而恭敬地以主教、教师和执事的身份行事，为天主的选民服务。他们中有些已安眠，有些仍存活。他们总是彼此和睦、相处平安、互相倾听。因此，他们完美地吻合在塔的建筑中。} \Quote{但那些从深处拖上来的石头，被放入建筑，并与先前已放在塔中的其余石头吻合的，是谁呢？} \Quote{他们是那些为主的缘故而受苦的人。} \Quote{但是，夫人，我想知道，那些从陆地上搬来的其他石头是谁？} \Quote{那些}，她说，\Quote{进入建筑而不加打磨的，是天主所赞许的，因为他们走在主的正道上，并奉行祂的诫命。} \Quote{但那些正被带来并放入建筑中的，是谁呢？} \Quote{他们是信德上年轻且忠信的人。他们受天使劝诫行善，因为在它们身上找不到任何邪恶。} \Quote{那么，那些被拒绝并抛弃的，是谁呢？} \Quote{这些是犯了罪并愿意悔改的人。为此，他们没有从塔中被远远扔出去，因为如果悔改，他们仍可在建筑中有用。那将要悔改的人，若悔改，将在信德上坚固——只要他们趁现在塔还在建造时悔改。因为如果建筑完工，就没有空间容纳任何人了，他必被抛弃。但这特权只属于那已被放在塔近处的人。}\par

% structure: p0034 source=h3 target=subsubsection reason=relative-heading-rank; base=h2

\subsubsection{第六章}

\Quote{至于那些被砍下并远远扔出塔外的，你想知道他们是谁吗？他们是邪恶之子，怀着伪善相信，邪恶从未离开他们。因此他们不能得救，因为出于他们的邪恶，他们不能被用于建筑。为此，他们被砍下并远远扔出，因主的震怒，他们激怒了祂。但我将向你解释你看见的其余那些大量躺着、未进入建筑的石头。粗糙的是那些认识了真理却未持守，也没有与圣人结合的人。因此他们不堪使用。} \Quote{那些有裂缝的是谁？} \Quote{这些是心中彼此不和、彼此不平安的人：他们表面彼此相安，但一分开，邪恶的念头便留在心中。这些，就是石头上的裂缝。而短一些的是那些确实相信、有较多正义，但也有不少邪恶的人，因此他们短而不全。} \Quote{但是，夫人，那些白色圆形、却未能嵌入塔的建筑中的，是谁呢？} 她回答说：\Quote{你还要愚笨无知多久，继续提出各种问题，却什么也不懂？这些是确有信德，但也有今世财富的人。所以，当患难来临，因着财富和生意，他们便否认了主。} 我回答她说：\Quote{那么，夫人，他们何时才能被用于建筑呢？} \Quote{当他们那如今诱惑他们的财富被削去时，他们就对天主有用。因为如同圆形石头若不被切去部分并抛弃便不能成为方形，同样，今世富有的人若不削减财富，就不能对主有用。先从你自身学这道理吧。你富有时，无用；但现在你有用且适于生命。要对天主有用；因为你也要被用作这些石头之一。}\par

% structure: p0036 source=h3 target=subsubsection reason=relative-heading-rank; base=h2

\subsubsection{第七章}

\Quote{现在你看见的其余远远扔出塔外、掉在公共道路上又滚入荒野之地的石头，是那些确实相信，却因怀疑而离开了正道的人。他们以为能找到更好的，便徘徊流浪，陷入悲惨，进入荒野之地。而那些掉进火里被烧毁的，是永远离开了永生天主的人；悔改的意念永未进入他们心中，因他们沉迷于情欲和所犯的罪行。你想知道其他那些落在水边却未能滚入水中的吗？这些是听到了圣言，并愿奉主名受洗的人；但当真理所要求的贞洁进入他们记忆时，他们便退缩，再次随从自己邪恶的欲望。} 她结束了关于塔的解释。但我仍不知羞耻，问她：\Quote{所有那些被抛弃、未能嵌入塔的建筑中的石头，还有可能悔改吗？他们在这塔中还有位置吗？} 她说：\Quote{悔改是可能的，但在这塔中他们找不到合适的位置。他们将被安放在另一个远为低劣的地方，而且只有当他们受折磨、度尽罪恶的日子之后才能如此。他们将因曾分享义德的圣言而得以转移。只有当心中兴起悔改所行之恶的意念时，他们才得脱离惩罚。但若这意念不进入他们心中，他们便不得拯救，因心硬。}\par

% structure: p0038 source=h3 target=subsubsection reason=relative-heading-rank; base=h2

\subsubsection{第八章}

当我停止就所有这些问题提问时，她对我说：\Quote{你愿意再看别的吗？} 我极其渴望看到更多，脸上洋溢着喜悦。她微笑着注视我，说：\Quote{你看见围绕塔楼的七位妇女吗？} \Quote{是的，夫人。} 我说。\Quote{这座塔楼}，她说，\Quote{是按照主的诫命由她们支撑的。现在请听她们的职责。第一位双手紧握的，称为信德。通过她，天主的选民得救。另一位卷起衣襟、行动有力的，称为节德。她是信德的女儿。凡跟随她的人，将在生命中幸福，因为他会克制自己远离一切恶行，相信若克制一切恶欲，他将继承永生。} \Quote{但是其他几位呢，} 我说，\Quote{啊，夫人，她们是谁？} 她对我说：\Quote{她们彼此是女儿。一位称为简朴，另一位称为纯朴，另一位称为贞洁，另一位称为明智，另一位称为爱德。当你实行她们母亲的一切工作时，你将能生活。} \Quote{我想知道，} 我说，\Quote{啊，夫人，每一位拥有怎样的能力。} \Quote{请听，} 她说，\Quote{她们拥有何种能力。她们的能力彼此制约，并按出生次序彼此跟随。因为从信德生出节德；从节德生出简朴；从简朴生出纯朴；从纯朴生出贞洁；从贞洁生出明智；从明智生出爱德。这些妇女的作为是纯洁、贞洁、神圣的。凡专心于这些作为并能持守它们的人，将与众位天主的圣人一同居住在塔楼内。} 然后我问及时期，是否现在已是终结。她大声呼喊说：\Quote{愚昧的人！你难道没看见塔楼仍在建造吗？当塔楼完工建成时，终结就来到；我向你保证，它很快就要完工。不要再问我问题。让你和所有的圣人满足于我所提醒你们的话，以及我更新你们心神的事吧。但请注意，这些启示不仅是为你们自己而赐下，更是为让你们能向众人显示。因为三天之后——你要留心记住这一点——我命令你把我将要对你说的一切话讲给圣人们听，使他们听了并实行，得以洗净自己的不义，你也同他们一起被洗净。}\par

% structure: p0040 source=h3 target=subsubsection reason=relative-heading-rank; base=h2

\subsubsection{第九章}

儿子们，请听我的话：我曾在极大的简朴、纯朴和贞洁中抚养了你们，这是由于主的仁慈，祂将祂的正义降在你们身上，使你们脱离一切不义和败坏，成为正义而圣洁的；但你们不愿停止不义。所以现在，请听我的话，彼此和睦，互相探访，互相承担重担，不要独自享用天主的受造物，而要慷慨分施给有需要的人。因为有些人因食物丰裕而使肉身虚弱，从而败坏了自己的肉身；而另一些没有食物的人的肉身却因营养不足而败坏。因此他们的身体日渐消瘦。这种饮食无度对你们这些丰裕而不分施给穷困者是有害的。要留意将来的审判。你们这些身居高位的人，趁着塔楼尚未完工，去寻找饥饿的人吧；因为塔楼完工后，你们想行善，却找不到机会了。所以，你们这些以财富为荣的人，要留心，免得穷困者叹息，他们的叹息上达于主，而你们连同所有财物被关在塔楼门外。因此现在我对你们这些管理教会、爱坐首位的人说：\Quote{不要像配药者一样。因为配药者把药品装在盒子里，而你们却把药品和毒药藏在心中。你们心硬，不愿洁净自己的心，也不愿将目标的一致与心的纯洁结合起来，以求得伟大君王的怜悯。孩子们，要当心，免得你们的这些纷争剥夺你们的生命。如果你们自己不受教导，怎能教导主的选民呢？所以你们要彼此教导，彼此和睦，好使我也能喜乐地站在你们的父面前，为你们众人向你们的主交账。}\par

% structure: p0042 source=h3 target=subsubsection reason=relative-heading-rank; base=h2

\subsubsection{第十章}

她停止对我说话后，那六位从事建筑的青年前来，将她送到塔楼那里，另外四个人抬起座椅，也搬到塔楼。后者的脸我没有看见，因为他们背对着我。当她离去时，我请求她向我启示她向我显现的三种形态的意义。她回答我说：\Quote{关于这些，你必须请另一个人向你揭示其意义。} 因为弟兄们，她在去年第一次神视中以一位极其年迈的老妇人形象向我显现，坐在椅子上。在第二次神视中，她的面容年轻，但皮肤和头发显出年迈，她站着对我说话。她也比第一次更喜乐。但在第三次神视中，她完全年轻，异常美丽，只是头发仍像老妇人；但她的面容洋溢喜乐，她坐在座位上。我对这些显现非常忧愁，因为我极想知道神视的含义。然后我在夜间的神视中看见那老妇人对我说：\Quote{每次祈祷都应伴随谦卑：所以你要斋戒，你必从主那里获得你所求的。} 于是我斋戒了一天。\par

就在那一夜，有一位青年向我显现，说：\Quote{你为什么屡次在祈祷中求启示？你要小心，免得因多求而损伤你的肉身；对这些启示应当知足。难道你还能看见比你已见的更大的启示吗？}我回答他说：\Quote{先生，我只求一件事，就是关于这三种形像的启示得以完全。}他回答我说：\Quote{你迟钝要到几时呢？但你的疑虑使你迟钝，因为你们的心没有转向主。}我却回答他说：\Quote{先生，我们要从你那里更准确地学习这些事。}\par

% structure: p0045 source=h3 target=subsubsection reason=relative-heading-rank; base=h2

\subsubsection{第十一章}

\Quote{那么，}他说，\Quote{关于你所询问的三种形像，请听。为什么在第一个异象中她以一个老妇人坐在椅子上的形像向你显现？因为你们的心神现在衰老枯干，因你们的软弱和疑虑而失去了力量。正如老年人不再指望恢复精力，只等最后的一眠，同样，你们因世俗的事务而衰弱，沉溺于懈怠，没有将你们的忧虑交托给主。因此你们的心神破碎，在忧愁中变得衰老。}\Quote{先生，那么我想知道她为什么坐在椅子上？}他回答说：\Quote{因为每个软弱的人因自己的软弱而坐在椅子上，好使他的软弱得到支撑。看，你已有了第一个异象的形像。}\par

% structure: p0047 source=h3 target=subsubsection reason=relative-heading-rank; base=h2

\subsubsection{第十二章}

\Quote{在第二个异象中，你看见她站着，面容年轻，比先前更喜乐；但她的皮肤和头发仍是老妇人的。请听，}他说，\Quote{这个比喻。当一个人变得有些年迈时，他因自己的软弱和贫困而对自己绝望，除了生命的最后一日之外别无所望。然后突然有一笔遗产留给他；他听说后，就起身，变得极其喜乐，穿上力量。如今他不再躺卧，而是站立；他先前因过去的作为已被摧毁的心神，得以更新，他不再坐着，而是精力充沛地行动。你们在听到天主所赐的启示时，也是如此。因为主怜悯了你们，更新了你们的心神，你们抛弃了软弱。力量在你们内兴起，你们在信德上变得坚强；主看见你们的力量，就喜乐。为此，祂向你们显示了塔的建造；并且祂还要向你们显示别的事，只要你们全心全意彼此和睦相处。}\par

% structure: p0049 source=h3 target=subsubsection reason=relative-heading-rank; base=h2

\subsubsection{第十三章}

\Quote{如今，在第三个异象中，你看见她更加年轻，高贵而喜乐，她的形貌美丽。因为，正如当一些好消息突然传给一个忧愁的人时，他立刻忘记先前的忧愁，只想着他所听到的好消息，从此为善而坚强，他的心神因所领受的喜乐而更新；同样，你们因看见这些美善之事而接受了心神的更新。至于你看见她坐在座位上，那意味着她的位置是稳固的，因为座位有四只脚，站立稳固。因为世界也是由四个元素维系在一起的。因此，那些完全并全心悔改的人，将变得年轻并牢固确立。如今你已完全得到了启示。不要再要求更多的启示。如果有任何事应当被启示，它将会启示给你。}\par

% structure: p0051 source=h2 target=subsection reason=relative-heading-rank; base=h2

\subsection{异象四}

\Emph{关于即将来临的考验与苦难。}\par

% structure: p0053 source=h3 target=subsubsection reason=relative-heading-rank; base=h2

\subsubsection{第一章}

前一个异象二十天后，我看见另一个异象，弟兄们——那即将来临的苦难的预象。我正沿着康帕尼亚路走向一座乡间别墅。那座房屋距公路大约十弗隆。那个地区人迹罕至。当我独自行走时，我祈祷主完成祂藉祂的圣教会所赐给我的启示，好使祂坚固我，并赐悔改给所有祂那些迷途的仆人，使祂伟大而荣耀的圣名得光荣，因为祂屈尊向我显示祂的奇迹。当我光荣祂、感谢祂时，有一个声音好像回答我说：\Quote{不要怀疑，赫尔玛；}我便开始自己思想，说：\Quote{我有什么理由怀疑呢？——我既已被主坚固，又看见了如此荣耀的景象？}我向前走了几步，弟兄们，看哪！我看见尘土甚至扬到天上。我对自己说：\Quote{是牲畜来了，扬起尘土吗？}那距离我大约一弗隆远。看哪！我看见尘土越来越浓，以致我想象那是从天主而来的什么东西。但太阳这时略微照耀出来，看哪！我看见一只像鲸鱼般的巨兽，从它口中飞出火蝗。那只兽的体长约有一百尺，它的头像一只瓮。我开始哭泣，呼求主救我脱离它。然后我记起我所听见的那句话：\Quote{不要怀疑，赫尔玛。}因此，我的弟兄们，既然怀着对主的信德，并记起祂所教导我的伟大事理，我便勇敢地面对那兽。那兽带着如此大的响声和力量而来，它本可毁掉一座城。我走近它，那巨兽伸展在地上，只现出它的舌头，丝毫不动，直到我走过它。那兽的头上又有四种颜色——黑色，然后是火红色和血色，然后是金色，最后是白色。\par

% structure: p0055 source=h3 target=subsubsection reason=relative-heading-rank; base=h2

\subsubsection{第二章}

我走过那野兽，又前进了大约三十英尺后，看，一位穿着如从洞房出来、全身雪白、脚穿白凉鞋、额头被面纱遮住、头戴风帽、头发雪白的处女迎向我。从先前的异象我认出这就是教会，便更加欢喜。她向我请安说：\Quote{愿你平安，人啊！}我回礼说：\Quote{夫人，愿你平安！}她回答我说：\Quote{有什么东西挡你的路吗？}我说：\Quote{我遇到一只如此巨大的野兽，足以毁灭民众，但靠着主的德能和祂的大仁慈，我逃脱了。}\Quote{你逃脱得好，}她说，\Quote{因为你将你的忧虑交给了天主，并向主敞开了你的心，相信除了祂伟大而荣耀的名号外，别无拯救。为此，主派了祂的天使，就是管理野兽的那位，名叫特格里，封住了它的口，使它不能撕裂你。你因着信德，又因面对如此巨兽没有怀疑，所以从大磨难中逃脱了。因此，你去向主的选民述说祂的奇事，告诉他们这野兽是预表将要来临的大磨难。如果你们预备自己，全心悔改，转向主，就有可能逃脱它，只要你们的心纯洁无瑕，并在余下的生命里无瑕地事奉主。将你们的忧虑交给主，祂必引导。你们这些怀疑的人，要信赖主，因为祂是全能的，能转离祂的忿怒，并将鞭笞降在怀疑者身上。那些听见这话却藐视的人有祸了：他们不如未曾出生。}\BibleRef{玛 26:24}\par

% structure: p0057 source=h3 target=subsubsection reason=relative-heading-rank; base=h2

\subsubsection{第三章}

我问她野兽头上的四种颜色。她回答我说：\Quote{你又在好奇这些事了。}\Quote{是的，夫人，}我说，\Quote{请告诉我它们是什么。}\Quote{听着，}她说：\Quote{黑色是我们居住的世界；火红色指出世界必因血与火而毁灭；金色是你们这些逃出世界的人。因为金子经火炼而成为有用，你们这些住在世上的人也如此受炼。因此，那些坚持到底、经受火炼的人必将因此被洁净。就如金子除去渣滓，你们也将除去一切忧愁和困苦，变得纯洁，适于建造塔楼。白色则是将要来临的世代，天主的选民将住在其中，因为天主拣选得永生的人将是无瑕而纯洁的。因此，不要止息地将这些话讲给圣徒们听。这就是将要来临的大磨难的预表。如果你们愿意，它就算不得什么。要记住先前所写的一切。}说着，她就离开了。但我没有看见她退到哪里去。然而有响声，我惊恐地转身，以为那野兽来了。\par

% structure: p0059 source=h2 target=subsection reason=relative-heading-rank; base=h2

\subsection{第五个异象}

\Emph{论诫命。}\par

我正在家中祈祷，坐在床榻上时，走进一位容貌光耀的人，像牧人一般穿着白山羊皮，肩上挎着行囊，手中拿着杖，向我请安。我回礼。他立刻坐在我旁边，对我说：\Quote{我是由一位极可敬的天使派来，与你同住你余生。}我以为他是来试探我的，便对他说：\Quote{你是谁？因为我知道我所受托付的那位。}他对我说：\Quote{你不认识我吗？}\Quote{不认识，}我说。\Quote{我，}他说，\Quote{就是你所受托付的那位牧者。}他说话间，容貌就改变了；于是我认出他就是我所受托付的那位。我立刻慌张起来，恐惧抓住我，因我如此卑鄙愚蠢地回答他而深深忧伤。但他回答我说：\Quote{不要慌张，倒要从我将给你的诫命中汲取力量。因为我被派来，}他说，\Quote{是要再次向你显示你先前所看见的一切，尤其是那些对你有用的。首先，写下我的诫命和比喻，其余的事我将依我所显示的写下。为此，}他说，\Quote{我命令你先写下诫命和比喻，以便你容易阅读，并能够遵守。}于是我就照他所吩咐的写下了诫命和比喻。因此，当你们听了这些，遵守并行走在其中，以纯洁的心践行，就必从主那里得到祂所应许给你们的一切。但如果你们听了却不悔改，反而继续加增你们的罪，那么你们就要从主那里得到相反的事。牧者，就是悔改的天使，命令我写下这一切话。\par

\SourceRecord{Translated by F. Crombie.；Ante-Nicene Fathers；Vol. 2.；Edited by Alexander Roberts, James Donaldson, and A. Cleveland Coxe.；Buffalo, NY: Christian Literature Publishing Co.,；1885.；<http://www.newadvent.org/fathers/02011.htm>.}

% structure: p0001 source=h1 target=section reason=page-title

\section{\WorkTitle{赫尔玛牧者（第二卷）}}

% structure: p0002 source=h2 target=subsection reason=relative-heading-rank; base=h2

\subsection{第一诫}

\Emph{论对天主的信德。}\par

首先，要相信有一位天主，祂创造了并完成了万有，使万有从无有而出。唯独祂能包含一切，但祂自己却不能被包含。因此，你要对祂怀有信德，并敬畏祂；敬畏祂，就要操练节制。持守这些诫命，你就能摆脱一切邪恶，并穿上义德的力量，并为天主而生活，只要你持守这条诫命。\par

% structure: p0005 source=h2 target=subsection reason=relative-heading-rank; base=h2

\subsection{第二诫}

\Emph{论避免恶言，并以纯朴施舍。}\par

祂对我说：\Quote{要纯朴无伪，你就会像那些不知邪恶毁灭人生命的孩童一样。首先，不要毁谤任何人，也不要乐意听他人毁谤。你若听了，若相信所听到的诽谤，就会分担毁谤者的罪；因为相信了，你也会说弟兄的坏话。这样，你便与毁谤者同负其罪。因为毁谤是邪恶的，是不安的魔鬼。它决不安居于和平，却总留在纷争中。你要谨防它，就能常与众人和睦。要穿上圣洁，其中没有邪恶的冒犯缘由，只有平稳与喜乐的行为。要行善；从你的劳苦所得、天主赐给你的酬报中，要纯朴地分施给一切贫乏者，不要犹豫该给谁、不该给谁。要给所有人，因为天主愿意祂的恩赐与众人分享。领受的人要向天主交代为什么领受、领受了什么。因为困苦者领受不会受责罚，但以虚假借口领受的，要受惩罚。施与的人则是无罪的。因为他既从主领受了，就纯朴地完成了他的服务，不犹豫该给谁、不该给谁。这样的服务，若以纯朴完成，在天主眼中是光荣的。因此，凡如此以纯朴服务的人，必为天主而生活。你要持守我给你的这些诫命，好使你和你一家的悔改得以在纯朴中被寻见，你的心也清洁无瑕。}\par

% structure: p0008 source=h2 target=subsection reason=relative-heading-rank; base=h2

\subsection{第三诫}

\Emph{论避免谎言，以及赫尔玛对其虚伪的悔改。}\par

祂又对我说：\Quote{要爱真理，让你的口只出真理，使天主放在你肉身内的灵，在众人面前显为真实的；那住在你内的主，也必受光荣，因为主在一切言语上都是真实的，祂内毫无虚假。因此，说谎的人否认主，并掠夺祂，不把所领受的归还给祂。因为他们从祂领受了一个无谬的灵。若他们把这灵不真实地还给祂，就玷污了主的诫命，成为掠夺者。}听了这些话，我痛哭流涕。祂见我哭泣，便对我说：\Quote{你为什么哭？}我回答：\Quote{先生，因为我不知道我是否能够得救。}\Quote{为什么？}祂说。我说：\Quote{因为先生，我一生从未说过真话，一向对所有人诡诈，对所有人都以谎言充当真理；从未有人反驳我，反而都相信我的话。我既如此行事，怎能活下去呢？}祂对我说：\Quote{你此刻的心情确是正确而健全的，因为你身为天主的仆人，本当在真理中行走，不可让邪恶的良心与真理之灵同流合污，也不可令圣洁真实的圣神忧伤。}我对祂说：\Quote{先生，我从未如此留意过这些话。}祂对我说：\Quote{现在你既听见了，就要持守它们，使得你从前在交易中所说的谎言，也能因你现今言语的真实而变得可信。因为即便那些谎言也能成为配得信任的。你若持守这些规诫，从今以后只说真理，你便能获得生命。凡听见这条诫命，并远离那极大的邪恶——谎言，就必为天主而生活。}\par

% structure: p0011 source=h2 target=subsection reason=relative-heading-rank; base=h2

\subsection{第四诫}

\Emph{论因奸淫休妻。}\par

% structure: p0013 source=h3 target=subsubsection reason=relative-heading-rank; base=h2

\subsubsection{第一章}

\Quote{我命令你，}他说，\Quote{要守护你的贞洁，不要让任何关于他人妻子、邪淫或类似不义的思想进入你的心；因为做这事你就犯大罪。但你若常思念自己的妻子，就永不会犯罪。因为这思想若进入你的心，你就会犯罪；同样，你若有其他恶念，你就犯罪。因为这思想在天主的仆人身上是大罪。但若有人行此恶事，便是为自己招致死亡。所以你要留心，远离这思想；因为纯洁所在之处，不义就不该进入义人的心。}我对他说：\Quote{先生，请允许我问您几个问题。}\Quote{说吧，}他说。我对他说：\Quote{先生，若有人有信主的妻子，若他发现她犯奸淫，他继续与她同居是否有罪？}他对我说：\Quote{只要他不知她的罪，丈夫与她同居不犯过。但若丈夫知道妻子失足，而妇人不悔改，反坚持她的邪淫，丈夫仍与她同居，他也有分于她的罪，同犯奸淫。}我对他说：\Quote{那么，先生，若妻子继续行恶，丈夫该做什么？}他说：\Quote{丈夫应休妻，独身自守。但若休妻而另娶，他也犯奸淫。}我对他说：\Quote{若被休的妇人悔改，愿回丈夫那里，她不应被丈夫重新接纳吗？}他对我说：\Quote{当然。若丈夫不接纳她，他就犯罪，并给自己带来大罪；因为他应当接纳悔改的罪人。但不可多次。因为天主的仆人只有一次悔改。因此，若被休的妻子可能悔改，丈夫在休妻后就不应另娶。在这事上，男女应同样对待。此外，不仅那些玷污自己肉身的人犯奸淫，那些在行为上仿效外邦人的也犯奸淫。因此，若有人坚持此类行为而不悔改，就离开他，停止与他同居，否则你就有分于他的罪。所以你们受命，应夫妇各自独居，因为这样的人可能悔改。但我并不，}他说，\Quote{给行这些事的机会，而是叫犯罪的人不再犯罪。至于他以往的过犯，天主能医治他，因为天主确实拥有万有之权柄。}\par

% structure: p0015 source=h3 target=subsubsection reason=relative-heading-rank; base=h2

\subsubsection{第二章}

我问他又说：\Quote{既然主恩允常与我同住，求您容忍我说几句话；因为我什么都不明白，我的心因我先前的生活方式变硬了。求您赐我理解，因为我极其迟钝，完全不懂。}他回答我说：\Quote{我主管悔改，并赐所有悔改的人理解。难道你不认为，}他说，\Quote{悔改是大智慧吗？因为悔改就是大智慧。那犯罪的人知道自己行恶于主前，记起自己所行的事，就悔改，不再行恶，反而慷慨行善，并因自己犯罪而谦卑折磨自己的灵魂。所以你看，悔改就是大智慧。}我对他说：\Quote{正因为如此，先生，我仔细查问一切，尤其因为我是一个罪人；为要知道该行什么善工，使我得以存活；因为我的罪又多又杂。}他对我说：\Quote{你若遵守我的诫命，并遵行它们，就必存活；凡听见并遵守这些诫命的，必生活于天主前。}\par

% structure: p0017 source=h3 target=subsubsection reason=relative-heading-rank; base=h2

\subsubsection{第三章}

我对他说：\Quote{我想继续提问。}\Quote{说吧，}他说。我说：\Quote{先生，我听见有些教师主张，除我们下到水里、领受赦免先前的罪的那次悔改外，没有别的悔改。}他对我说：\Quote{你听到的是正确的道理；因为事实确实如此。因为那领受赦罪的人，不应再犯罪，而应生活在纯洁中。既然你殷勤查问一切，我也将这事指示你，但不是给那些将要相信或最近相信主的人提供犯错误的机会。因为现在已相信和将要相信的人，没有悔改的机会（为他们的罪）；他们只有赦免先前的罪。因为在那些日子之前蒙召的人，主为他们设立了悔改。因为主认识人心，预知万事，知道人的软弱和魔鬼的多种诡计，魔鬼将加害天主的仆人并恶待他们。因此，慈悲的主怜悯了祂手的工作，为他们设立了悔改，并将这悔改的权柄托付给我。所以我告诉你，若有人受魔鬼诱惑，在那伟大而神圣的召叫中主召选祂的子民得永生之后犯了罪，他只有一次悔改的机会。但若此后常犯罪，然后又悔改，此人的悔改将无济于事，因为他难以得生。}我说：\Quote{先生，我觉得生命已回到我身上，因留心听这些诫命；因为我知道，若我今后不再犯罪，我必得救。}他说：\Quote{你必得救，你和所有遵守这些诫命的人。}\par

% structure: p0019 source=h3 target=subsubsection reason=relative-heading-rank; base=h2

\subsubsection{第四章}

我又问他说：\Quote{先生，因你如此耐心听我，可否也给我显示这事？}他说：\Quote{说吧。}我便说：\Quote{若丈夫或妻子死了，鳏夫或寡妇再婚，他或她是否犯罪？}他说：\Quote{再婚并无罪；但若他们不续弦，在主前获得更大的尊荣和荣耀；但若他们结婚，也不犯罪。因此，务要保守你的贞洁和纯洁，你将因天主而活。我现在赐给你的诫命，以及我将要赐给你的，从今以后，即从你交托给我的那一天起，你要遵守，我就住在你家中。你若遵守我的诫命，你从前的罪必得赦免。凡遵守我这些诫命，并行走在这贞洁中的人，都必得赦免。}\par

% structure: p0021 source=h2 target=subsection reason=relative-heading-rank; base=h2

\subsection{第五诫}

\Emph{论心中的忧苦与忍耐}\par

% structure: p0023 source=h3 target=subsubsection reason=relative-heading-rank; base=h2

\subsubsection{第一章}

\Quote{你要忍耐，}他说，\Quote{并要有明智，这样你必管辖一切邪恶的行为，并成就一切义德。因为你若忍耐，住在你内的圣神必是纯洁的。祂不会被任何邪灵昏暗，反而住在宽阔之地，喜乐欢欣；并同祂所住的器皿一同喜乐地事奉天主，内心大有平安。但一旦有怒气爆发，温柔的圣神立即受困，因没有纯洁之处，祂便寻求离去。因为祂被卑劣的灵窒息，不能按祂所愿事奉主，因为忿怒玷污了祂。主住在宽容中，魔鬼却住在忿怒中。故此，两个灵同住一室时，彼此不和，使那住有它们的人烦恼。因为若将极小的一块苦艾放入一罐蜂蜜中，岂非整个蜂蜜被败坏，那极小的一块苦艾岂非完全夺去蜂蜜的甘甜，以致它不再给主人任何愉悦，反而变苦，失去用处？但若不将苦艾放入蜂蜜，蜂蜜便保持甘甜，对主人有用。你看，忍耐比蜂蜜更甘甜，对天主有用，主住在其中。忿怒却是苦而無用的。若将忿怒掺入忍耐，忍耐便被玷污，它的祈祷对天主便不再有用。}我说：\Quote{先生，我愿知道忿怒的威力，好能防备它。}他说：\Quote{你若不对它设防，你和你的家便失去一切得救的希望。因此，你要防备它。因为我与你同在，凡全心悔改的都必远离它。因为我必与他们同在，且要拯救他们众人。因为众人都藉至圣的天使得以成义。}\par

% structure: p0025 source=h3 target=subsubsection reason=relative-heading-rank; base=h2

\subsubsection{第二章}

\Quote{现在请听，}他说，\Quote{忿怒的行为何等邪恶，它如何藉其行为推翻天主的仆人，使他们离开正义。但它不能推翻那些充满信德的人，也不能在他们身上行动，因为主的大能与他们同在。它推翻的是那些轻率多疑的人。因为一旦它看见这样的人站立坚定，它就投入他们心中，为了一件小事，男人或女人就因日常生活中的事端而变苦，例如为食物，或某句多余的话，或为某个朋友，或某项赠礼、债务，或诸如此类愚昧的事。因为这一切对天主的仆人而言，都是愚妄、空虚、无益的。忍耐却是伟大、有力、强壮、在宽阔中平静、喜乐、欢欣、无忧无虑、时时光荣天主、心中毫无苦毒、持续温柔安静。忍耐与那些有完备信德的人同在。忿怒却是愚拙、善变、无知的。愚拙生出苦毒，苦毒生出忿怒，忿怒生出疯狂。这疯狂由如此多的邪恶而生，终致大而难恕的罪。因为当所有这些灵同住在一器皿中，而圣神也住在其中时，这器皿不能容纳它们，便要溢出。温柔之神不惯与邪灵同住，也不惯与刚硬同住，便离开这样的人，寻求与温柔和平安同住。当祂离开所住的人时，那人便空虚了正义之神；从此充满邪灵，在一切行动中无法无天，被邪灵忽东忽西地牵引，心中对一切善事全然黑暗。凡忿怒的人都遭遇如此。因此，你要远离那极邪恶的忿怒之灵，穿上忍耐，抵抗忿怒和苦毒，你便与主所喜爱的纯洁同在。你要留意，切莫疏忽了这诫命；因为你若遵守这诫命，就能遵守我将赐给你的其他一切诫命。因此，你要在这些诫命上刚强，穿上能力；凡愿意行走在其中的，都要穿上能力。}\par

% structure: p0027 source=h2 target=subsection reason=relative-heading-rank; base=h2

\subsection{第六诫}

\Emph{如何识别每人身上的两个灵，及如何分辨一方与另一方的建议。}\par

% structure: p0029 source=h3 target=subsubsection reason=relative-heading-rank; base=h2

\subsubsection{第一章}

他说：\Quote{我在第一条诫命里给了你指示，要你关心信德、敬畏和自制。}我说：\Quote{是的，先生。}他说：\Quote{现在我想向你展示这些德能的能力，好叫你明白每种能力拥有什么力量。因为它们是双重的，而且同等关乎义人和不义的人。因此，你要信赖义人，但不要信任不义的人。因为义路是直的，但不义之路是弯曲的。但要走在笔直平坦的路上，不要理会弯曲的路。因为弯曲的路没有道路，却有许多无路之处和绊脚石，且崎岖多荆棘。它伤害行走其上的人。但行在直路上的人，行走平稳，不磕绊，因为既不崎岖也无荆棘。你看到了，走这条路更美。}我说：\Quote{我希望走这条路。}他说：\Quote{你会走这条路；凡全心归向主的人，必行走在其中。}\par

% structure: p0031 source=h3 target=subsubsection reason=relative-heading-rank; base=h2

\subsubsection{第二章}

他说：\Quote{现在请听关于信德的事。有两位天使与一个人同在——一位是义的天使，另一位是邪恶的天使。}我对他说：\Quote{先生，我如何能知道这些天使的能力，因为两位天使都住在我里面？}他说：\Quote{请听，并要明白它们。义的天使是温良、谦逊、柔和、平安的。因此，当他进入你的心时，他立刻与你谈论正义、纯洁、贞洁、知足，以及每一件正义的行为和光荣的德行。当这一切进入你的心时，你就知道义的天使与你同在。这些是义的天使的作为。那么，你要信赖他和他的工作。现在看看邪恶天使的作为。首先，他是易怒、苦毒、愚昧的，他的作为是邪恶的，败坏天主的仆人。因此，当他进入你的心时，你要从他的作为认出他。}我对他说：\Quote{先生，我如何识别他，我不知道。}他说：\Quote{请听并要明白。当愤怒或苛责临到你时，你就知道他在你里面；当你被许多事务的渴望、最丰盛的佳肴、醉酒的狂欢、各种奢侈、不当之事、对女人的贪恋、欺诈、骄傲、喧嚷以及类似的事攻击时，你也知道是这样。当这些进入你的心时，你就知道邪恶的天使在你里面。既然你知道他的作为，就离开他，在各方面不要信任他，因为他的行为是邪恶的，对天主的仆人无益。这样，两位天使的行动你们就知道了。要明白它们，信赖义的天使；但离开邪恶的天使，因为他的教导在每件事上都是邪恶的。因为即使一个人极其忠诚，但这天使的思想进入他的心，那人或女人必定犯罪。反之，即使一个人或女人极其邪恶，但如果义的天使的作为进入他的心或她的心，他或她必定行善。因此你看，跟随义的天使是好的，而告别邪恶的天使是好的。}\par

\Quote{这条诫命展示了信德的作为，好叫你信赖义天使的作为，并因遵行它们而活于天主。但要相信邪恶天使的作为是艰难的。如果你拒绝做它们，你必活于天主。}\par

% structure: p0034 source=h2 target=subsection reason=relative-heading-rank; base=h2

\subsection{第七诫命}

\Emph{论敬畏天主，不畏魔鬼。}\par

他说：\Quote{敬畏主，遵守祂的诫命。\BibleRef{训 12:13} 因为如果你遵守天主的诫命，你必在每一行动上有能力，每一个你的行动都将无可比拟。因为敬畏主，你必善作一切事。这就是你应当具有的敬畏，好使你得救。但不要畏惧魔鬼；因为敬畏主，你必制胜魔鬼，因为他里面没有能力。但在他里面没有能力的人，绝不应成为畏惧的对象；但在他里面有荣耀能力的那位，才是应当畏惧的。因为凡有能力的人都应当被畏惧；但无能的人被众人藐视。因此，要畏惧魔鬼的作为，因为它们是邪恶的。因为敬畏主，你就不会做这些作为，反而会戒避它们。因为敬畏有两种：如果你不愿意做恶事，敬畏主，你就不会做；但另一方面，如果你愿意做善事，敬畏主，你就会去做。因此，主的敬畏是强大、伟大而荣耀的。你们要敬畏主，就必活于祂；凡敬畏祂并遵守祂诫命的人，必活于天主。}我说：\Quote{先生，为什么关于那些遵守祂诫命的人，你说他们必活于天主？}他说：\Quote{因为一切受造物都敬畏主，但并非一切受造物都遵守祂的诫命。只有那些敬畏主并遵守祂诫命的人，才有在天主内的生命；但那些不遵守祂诫命的人，在他们里面没有生命。}\par

% structure: p0037 source=h2 target=subsection reason=relative-heading-rank; base=h2

\subsection{第八诫命}

\Emph{我们应该避开恶事，行善事。}\par

\Quote{我告诉过你，}他说，\Quote{天主的受造物是双重的，因为节制也是双重的；因为有些情况需要节制，有些情况则不需要节制。}\Quote{先生，请告诉我，}我说，\Quote{在哪些情况下需要节制，在哪些情况下不需要。}\Quote{在恶事上节制自己，不要去做；但在善事上不要节制，而要去做。因为如果你在行善时节制自己，你就会犯大罪；但如果你节制自己不去做恶事，你就是在实行伟大的义德。因此，你要戒绝一切不义，只行善事。}\Quote{先生，}我说，\Quote{我们必须节制自己的恶行有哪些？}\Quote{听吧，}他说：\Quote{就是奸淫、淫乱、非法狂欢、邪恶奢华、贪食美味和财富挥霍、自夸、傲慢、无礼、谎言、诽谤、虚伪、记仇和一切毁谤。这些是人生中最邪恶的行为。因此，天主的仆人要节制自己，不做这些行为。因为凡不节制这些的，不能为天主而活。那么，再听与这些相伴的行为。}\Quote{先生，}我说，\Quote{还有其他恶行吗？}\Quote{有，}他说，\Quote{有很多，天主的仆人也必须节制自己——偷窃、说谎、抢劫、假见证、欺诈、邪恶情欲、诡诈、虚荣、自夸，以及其他类似的一切恶习。}\Quote{你不认为这些确实邪恶吗？}\Quote{对天主的仆人来说极其邪恶。天主的仆人必须节制所有这一切。因此，你要节制自己远离这一切，好能为天主而活，你将被列在那些在这些事上节制自己的人中。那么，这些就是你必须节制的事。}\par

\Quote{但是请听，}他说，\Quote{那些你不必节制、反而应当去做的事。在善事上不要节制自己，而要去做。}\Quote{先生，请告诉我，}我说，\Quote{善行的本质是什么，好使我行走在其中、谨守它们，这样我就能得救。}\Quote{听着，}他说，\Quote{那些你应当做、且无需节制的善行。首先有信德，然后是敬畏主、爱德、和睦、正义的话语、真理、忍耐。在人的生命中，没有比这些更好的了。如果有人遵行这些，且不节制自己，他的一生就有福了。此外，还有这些随附的善行：帮助寡妇、照料孤儿和穷人、从困境中解救天主的仆人、款待旅客——因为款待是行善的园地——从不与人作对、安静、比众人需求更少、尊敬长者、实行义德、照顾弟兄姊妹、忍受侮辱、长期忍耐、鼓励心灵软弱的人、不把那些因犯罪而跌倒的人从信仰中抛弃，而是使他们回转、恢复心灵的平安、劝诫罪人、不压迫负债人和穷人，以及任何类似的行为。你看这些是善行吗？}他说。\Quote{先生，有什么比这些更好呢？}我说。\Quote{那么行走在其中吧，}他说，\Quote{不要自我节制，远离它们，你将为天主而活。因此，遵守这条诫命。你若行善，且不节制自己，你将为天主而活。所有这样行的人，都将为天主而活。同样，你若拒绝作恶，并且节制自己，你将为天主而活。所有遵守这些诫命、行走在其中的人，都将为天主而活。}\par

% structure: p0041 source=h2 target=subsection reason=relative-heading-rank; base=h2

\subsection{第九诫}

\Emph{应不断以坚定不移的信心向天主祈祷。}\par

祂对我说：\Quote{你要从自己身上除去疑虑，不要犹豫向主祈求，不要在心里说：\InnerQuote{我怎能向主祈求并从祂领受呢？看，我犯了如此多的罪得罪了祂。}不要这样与自己理论，但要全心归向主，毫不疑虑地向祂祈求，你便会知道祂的仁慈何其丰厚；祂绝不会舍弃你，反而要满全你灵魂的请求。因为祂不像世人那样记恨别人对他们的恶行；但祂自己却不记恶，且怜悯祂自己的受造物。因此，你要从心里清除这世界的一切虚荣，以及前面所提及的话语，并向主祈求，你必会领受一切，凡你毫无疑虑地向主所提出的请求，没有一样会被拒绝。但如果你心中疑虑，你就不会得到任何请求。因为那些对天主疑虑的人，是\OriginalTerm{双魂}{double-souled}的人，他们的请求一样也得不到。但那在信德上成全的人，凡事信赖主，就必获得；因为他们祈求时毫无疑虑，也不是双魂的人。因为每一个双魂的人，即使他悔改，也难以得救。因此，你要从心里清除一切疑虑，穿上信德，因为信德是坚强的；要信赖天主，你必从祂那里获得你所求的一切。如果你在向主祈求之后，得到所求的比你所期望的更迟，不要因为你没有很快得到你灵魂的请求而疑虑；因为通常是由于某种试探或某种你所不知道的罪，你才迟迟得不到所求的。所以，不要停止你灵魂的请求，你必会得到。但如果你在请求中变得疲倦和动摇，你就要责怪自己，而不是那不给你施予的主。你要考虑这种疑虑的心态，因为它是邪恶而愚昧的，并使许多人完全背离信德，即使他们原本非常坚强。因为这疑虑是魔鬼的女儿，并对天主的仆人行极其邪恶的事。因此，你要藐视疑虑，在一切事上胜过它；要穿上信德，因为信德是坚强有力的。因为信德应许一切，成就一切；但疑虑本身没有彻底的信德，在它所从事的一切工作中都失败。你看到了，}祂说，\Quote{信德是从上而来的——来自主——并有极大的能力；但疑虑是属地的精神，来自魔鬼，毫无能力。因此，你要事奉那有能力的东西，即信德，并远离那无能力的疑虑，这样你便能为天主而活。凡心思专注在这些事上的人，都将为天主而活。}\par

% structure: p0044 source=h2 target=subsection reason=relative-heading-rank; base=h2

\subsection{训言第十}

\Emph{论忧愁，以及不令在我们内的天主之圣神忧愁。}\par

% structure: p0046 source=h3 target=subsubsection reason=relative-heading-rank; base=h2

\subsubsection{第一章}

\Quote{你要从你身上除去，}祂说，\Quote{忧愁；因为她是疑虑和愤怒的姊妹。}\Quote{主，怎么，}我说，\Quote{她是这些的姊妹呢？因为愤怒、疑虑和忧愁彼此似乎很不相同。}\Quote{愚昧的人啊，你没有察觉忧愁比所有其他的灵更邪恶，对天主的仆人最可怕，比所有其他的灵更能摧毁人，并能压灭圣神，但另一方面又能拯救人吗？}\Quote{主，我是愚昧的，}我说，\Quote{不明白这些比喻。因为我不明白她如何既能压灭又能拯救。}\Quote{你听，}祂说。\Quote{那些从未寻求过真理，也未探求过\OriginalTerm{天主性}{Divinity}的本性，而只是单纯相信的人，当他们投身于并陷入事业、财富、外邦人的友谊以及许多这世界的其他行为时，就不能领悟天主性的比喻；因为他们的心思被这些行为蒙蔽，他们被败坏而变得枯干。正如美丽的葡萄树，若被忽视，就会被荆棘和各种植物所枯萎；同样，那些相信了而后来陷入上述许多行为的人，他们的心思就迷失了，对正义完全失去了理解力；因为他们即使听到正义，他们的心思也被事业占据，一点也不留心。另一方面，那些有敬畏天主之心，并寻求\OriginalTerm{天主性}{Godhead}和真理，且他们的心转向主的人，就很快领悟并理解对他们所说的话，因为他们心中有对主的敬畏。因为哪里主居住，哪里就有许多理解。因此，你要紧靠主，就必能领悟并通晓一切。}\par

% structure: p0048 source=h3 target=subsubsection reason=relative-heading-rank; base=h2

\subsubsection{第二章}

\Quote{那么，请听，}他说，\Quote{愚昧的人，忧愁如何压制并驱出圣神，另一方面又如何拯救。疑惑的人尝试做任何事，却因他的疑惑而失败时，这忧愁就进入那人内，并使圣神忧愁，并将祂压制出去。然后，另一方面，当愤怒附着于一个人，涉及某事，而他感到痛苦时，忧愁就进入那被激怒的人的心中，他因所做的事而忧愁，并悔改自己行了恶事。那么，这忧愁似乎伴随着救恩，因为那人在行了恶事后悔改了。这两种行动都使圣神忧愁：疑惑，因为它没有达成目标；愤怒使圣神忧愁，因为它行了恶事。这两者——疑惑和愤怒——都是令圣神忧愁的。因此，你要从你身上除去忧愁，不要压制住在你内的圣神，免得祂向你求告天主，并离开你。因为天主赐予我们住在这身体内的圣神，不能忍受忧愁或困迫。因此，你要穿上喜乐，这喜乐总是令天主喜悦和悦纳，并要在其中喜乐。因为每一个喜乐的人都行善事，思念善事，轻视忧愁；但忧愁的人总是行为邪恶。首先，他行为邪恶是因为他使圣神忧愁，圣神原是赐给人作为喜乐之神的。其次，使圣神忧愁，他行事不义，既不向主祈求，也不向祂认罪。因为忧愁的人的祈求没有力量上升到天主的祭台。}\Quote{为什么？}我说，\Quote{忧愁的人的祈求不能上升到祭台吗？}\Quote{因为，}他说，\Quote{忧愁坐在他心里。因此，忧愁与他的祈求混合，不容许祈求纯洁地上升到天主的祭台。因为正如醋和酒混合在同一器皿中，不能产生同样的愉悦（如纯酒单独产生的），同样，忧愁与圣神混合也不能产生同样的祈求（如单独由圣神产生的）。你要洁净自己脱离这邪恶的忧愁，你就将为天主而生活；所有驱除忧愁、穿上完全喜乐的人，都将为天主而生活。}\par

% structure: p0050 source=h2 target=subsection reason=relative-heading-rank; base=h2

\subsection{第十一诫命}

\Emph{圣神与先知应由其行为辨别；亦论两种灵。}\par

祂指给我看一些人坐在长凳上，另一个人坐在椅子上。祂对我说：\Quote{你看见坐在长凳上的人吗？} \Quote{我看见，先生，}我说。\Quote{这些人，}祂说，\Quote{是有信德的人，而坐在椅子上的那人是假先知，败坏天主的仆人的心思。他所败坏的，不是有信德的人，而是那些心怀疑虑的人。这些心怀疑虑的人便到他那里，如同到占卜者那里一般，询问他们将会发生什么事；而他，这假先知，本身没有天主圣神的能力，就按他们的询问，按他们的邪恶欲望回答他们，并按照他们自己的愿望使他们的灵魂充满期待。因为他自己本是空虚的，就给空虚的询问者空虚的回答；因为每一个回答都是针对人的空虚而发的。他有时也说出一些真话；因为魔鬼用他自己的灵充满他，希望他能战胜某些义人。因此，凡是那些在主内信德坚固、并披戴真理的人，不与这样的灵来往，反而远离它们；但那些心志不定、常常悔改的人，却去寻求占卜，如同外邦人一般，并因他们的偶像崇拜而给自己带来更大的罪。因为就任何行为向假先知询问的人，是拜偶像的人，没有真理，且是愚昧的。因为天主所赐的灵，不需要被人询问；这样的灵既有天主的能力，就自然说出一切，因为它来自上界，来自天主圣神的能力。但那受了询问、按人的欲望而发言的灵，是属地的、轻浮的、无力的，并且若不被询问，它就完全沉默。} \Quote{那么，先生，}我说，\Quote{一个人如何知道他们中哪一个是先知，哪一个是假先知呢？} \Quote{我要告诉你，}祂说，\Quote{关于这两种先知的情形，然后你就可以按照我的指示辨别真先知和假先知。你要通过他的生活来试验那拥有天主圣神的人。首先，那拥有自上天而来的天主圣神的人，是温良的、和平的、谦卑的，远离一切不义和这世界的虚妄欲望，比众人更少需求，并且当他被问时，他不做回答；他也不私下说话，也不在人希望那灵说话时让圣神说话，而只在天主愿意它说话时，它才说话。因此，当一个拥有天主圣神的人来到一群信德坚固、对天主圣神有信德的义人当中，这群人向天主祈祷时，那与他相配的先知之神的使者就充满他；那人被圣神充满，就按主的意愿向群众说话。这样，天主之神就显明出来。因此，凡从天主之神而来的能力，都属于主。那么，请听，}祂说，\Quote{关于那属地的、空虚的、无力的、愚昧的灵。首先，那似乎拥有这灵的人自高自大，想要居首位，并且胆大、厚颜、多言，生活在许多奢靡和许多其他迷惑之中，并为他的预言收取酬报；如果他得不到酬报，他就不预言。那么，天主圣神能收取酬报并预言吗？天主的先知不可能这样做，但这类先知是被属地的灵所附的。其次，它从不接近义人的集会，反而避开他们。它和心怀疑虑的人及空虚的人交往，在角落里向他们预言，欺骗他们，按他们的欲望说些空洞的话：因为它所回答的那些人本是空虚的。因为空虚的器皿放在空虚的器皿中，不会被压碎，而是彼此相符。因此，当它来到一群拥有天主之神的义人当中，他们祈祷时，那人就变得空虚，属地的灵因恐惧而从他身上逃离，那人就变哑，完全被压碎，不能说话。因为如果你把一仓库装满酒或油，在酒或油的器皿中间放一个空罐子，你若想清理仓库，就会发现那罐子仍像你放进去时一样是空的。同样，空虚无物的先知来到义人的灵中时，[离开时]仍像他们来时一样。这就是两种先知的生活方式。你要通过那自称受灵感之人的作为和生活来试验他。至于你，要相信那来自天主、有能力的灵；但属地的、空虚的灵，一点也不要相信，因为它里面没有能力：它来自魔鬼。那么，请听我将要对你说比喻。拿一块石头，扔到天上，看你能不能碰到它。或者，再拿一注水，喷向天上，看你能不能穿透天空。} \Quote{先生，}我说，\Quote{这些事怎么能发生呢？因为两者都是不可能的。} \Quote{正如这些事，}祂说，\Quote{是不可能的，同样，属地的灵也是无力和无能的。但另一方面，请看来自上界的能力。冰雹小得像一粒微小的谷粒，但落在人头上时，给人带来多大的烦恼！或者，再取一滴从水罐滴落在地上的水滴，却能穿透石头。你看到，来自上界的最微小的事物，落在地上时也具有大能。同样，从天主之神，即来自上界的，也是有能力的。所以，要相信那一个神，但不要与另一个神有任何瓜葛。}\par

% structure: p0053 source=h2 target=subsection reason=relative-heading-rank; base=h2

\subsection{第十二条诫命}

\Emph{论双重欲望。天主的诫命可以遵守，信徒不应惧怕魔鬼。}\par

% structure: p0055 source=h3 target=subsubsection reason=relative-heading-rank; base=h2

\subsubsection{第一章}

他對我說：\Quote{應棄絕一切邪惡的慾望，披上良善而純潔的慾望；因為披上這慾望，你必憎恨邪惡的慾望，並能隨意克制自己。因為邪惡的慾望是野蠻的，難以馴服。它極其可怕，因它的野蠻而極度吞噬人。尤其天主的僕人若陷在其中又缺乏理解，必被它可怕地吞噬。此外，它吞噬所有沒有披上良善慾望之衣、而糾纏混雜於此世的人，將這些人交付死亡。}\Quote{那麼，先生，}我說，\Quote{邪惡慾望的哪些行為把人交付死亡？請讓我知道，我好遠離它們。}\Quote{那麼，請聽邪惡慾望殺害天主僕人的作為。}\par

% structure: p0057 source=h3 target=subsubsection reason=relative-heading-rank; base=h2

\subsubsection{第二章}

\Quote{首先，是貪戀他人妻子或丈夫的慾望，以及放縱、許多無用的美食與飲料，和許多其他愚昧的奢華；因為在天主僕人中，一切奢華都是愚昧而空虛的。這些就是殺害天主僕人的邪惡慾望。因為這邪惡的慾望是魔鬼的女兒。你必須遠離邪惡的慾望，好能因遠離而為天主生活。但那些被它們制服且不抵抗的人，最終必滅亡，因為這些慾望是致命的。所以，你當披上正義的慾望，並以敬畏主為武器，抵抗它們。因為敬畏主居住在良善的慾望中。但若邪惡的慾望見你以敬畏天主為武裝並抵抗它，它就必遠離你，不再向你顯現，因為它懼怕你的武器。那麼，帶著你戰勝它所得的花冠，前往正義的慾望那裡，將你所獲的獎賞獻給它，並按它的意願服事它。你若服事良善的慾望並順從它，就必制服邪惡的慾望，使它隨你的意願服從你。}\par

% structure: p0059 source=h3 target=subsubsection reason=relative-heading-rank; base=h2

\subsubsection{第三章}

\Quote{我想知道，}我說，\Quote{我該如何服事良善的慾望。}\Quote{請聽，}他說：\Quote{你將實踐正義與德行、真理與敬畏主、信德與溫和，以及諸如此類的一切美善。實踐這些，你就成為天主所喜悅的僕人，並為祂而生活；凡服事良善慾望的人，都將為天主而生活。}\par

他結束了十二條誡命，對我說：\Quote{現在你有了這些誡命。按此而行，並勸勉你的聽眾，使他們的悔改在餘生中保持純潔。忠實履行我現今交託給你的這項職務，你將成就許多事。因為你將在那些將要悔改的人中獲得恩寵，他們必留意你的話；因為我與你同在，並將強迫他們服從你。}我對他說：\Quote{先生，這些誡命偉大、良善、光榮，能使能實行的人心中喜樂。但我不知道人能否遵守這些誡命，因為它們極其困難。}他回答我說：\Quote{你若確信它們能被遵守，你就容易遵守，它們並不困難。但你若想像它們不能被遵守，那麼你就不會遵守。我現在告訴你：你若不去遵守，反輕忽它們，你必不得救，你的兒女、你的家也不得救，因為你已斷定這些誡命人不能遵守。}\par

% structure: p0062 source=h3 target=subsubsection reason=relative-heading-rank; base=h2

\subsubsection{第四章}

這些話他以極其憤怒的語氣對我說，使我驚慌，對他極其恐懼，因為他的容貌改變，人無法承受他的怒氣。但見我完全動搖慌亂，他便開始以更溫和的語氣對我說：\Quote{愚昧、無知又遲疑的人哪！你豈不曉得天主的榮耀何其偉大、剛強而奇妙？祂為人創造了世界，使萬物服從於他，賜他權柄管理天下一切。那麼，人既是天主造物的主宰，管理萬有，難道不能也成為這些誡命的主嗎？}他說：\Quote{因為那心中有主的人，也能成為一切及每一條誡命的主。但那些只在口唇上有主\BibleRef{依 29:13}、內心卻剛硬\BibleRef{若 12:40}、遠離主的人，誡命對他們是困難而艱澀的。所以，你這個在信德上空虛搖擺的人，將主置於心中，你就會知道沒有比這些誡命更容易、更甘飴、更易管理的了。回轉吧，你們這些行走在魔鬼誡命中、處於艱苦、苦澀而野蠻的放縱中的人，不要懼怕魔鬼；因為牠在你身上毫無能力，因為我——悔改的天使，是管轄牠的——與你同在。魔鬼只有恐懼，但牠的恐懼沒有力量。所以不要怕牠，牠必逃離你。}\par

% structure: p0064 source=h3 target=subsubsection reason=relative-heading-rank; base=h2

\subsubsection{第五章}

我对他说：\Quote{先生，请听我说片刻。}他说：\Quote{请说。}我说：\Quote{先生，人急切地想遵守天主的诫命，没有一个人不向主祈求，为这些诫命赐予力量，好能服从它们；但魔鬼强硬，控制着他们。}他说：\Quote{他不能控制天主的仆人，他们全心寄望于祂。魔鬼能与他们搏斗，却不能推翻他们。因此，你若抵抗他，他必被制伏，羞愧地逃离你。}他说：\Quote{因此，许多空的人害怕魔鬼，认为他有能力。当一个人用美酒装满很好的坛子，其中几坛空着，他就来到坛前，不看满的坛，因为他知道它们是满的；却看空的，怕它们变酸。因为空坛很快变酸，酒的美味就消失了。魔鬼同样去到所有天主的仆人那里试探他们。那些信德充满的人，就强烈地抵抗他，他便离开他们，无路可入。然后他到空的里面，找到进入之路，就在他们身上造出他所愿意的一切，他们就成了他的仆人。}\par

% structure: p0066 source=h3 target=subsubsection reason=relative-heading-rank; base=h2

\subsubsection{第六章}

他说：\Quote{但我，悔改的天使，对你们说：不要害怕魔鬼；因为我被派遣，} \Quote{与你们这些全心悔改的人同在，并使你们在信德上坚强。那么，你们这些因自己的罪而对生命绝望的人，又增添罪恶、加重生命重担的人，要信赖天主；因为如果你们全心归向主，在余下的日子里实践正义，按祂的旨意事奉祂，祂必医治你们从前的罪，你们将有能力控制魔鬼的工作。至于魔鬼的威胁，一点也不要怕，因为他像死人筋骨一样无力。所以，听从我，敬畏那有能力拯救和毁灭的那位，\BibleRef{玛 10:28}，并遵守祂的诫命，你们将活于天主。}我对他说：\Quote{先生，我现在在主的一切规诫上坚强起来，因为你与我同在；我知道你将粉碎魔鬼的一切权势，我们将统治他，胜过他的一切工作。我也希望，先生，靠主加强我，能遵守你吩咐我的这一切诫命。}他说：\Quote{你必遵守，} \Quote{如果你的心在主面前纯洁；凡洁净自己心灵脱离此世虚妄欲望的人，也必遵守，他们将活于天主。}\par

\SourceRecord{Translated by F. Crombie.；Ante-Nicene Fathers；Vol. 2.；Edited by Alexander Roberts, James Donaldson, and A. Cleveland Coxe.；Buffalo, NY: Christian Literature Publishing Co.,；1885.；<http://www.newadvent.org/fathers/02012.htm>.}

% structure: p0001 source=h1 target=section reason=page-title

\section{\WorkTitle{黑马牧人书}（第三卷）}

% structure: p0002 source=h2 target=subsection reason=relative-heading-rank; base=h2

\subsection{比喻一}

\Emph{在此世我们没有永存的城，我们应当寻求来世之城。}\par

\Quote{你们知道：你们这些天主的仆人居住在异乡；因为你们的城离此城很远。如果这样，}他继续说，\Quote{你们知道你们将要居住的城，为什么在此地购置田地、作昂贵的准备、积累住所和无用的建筑呢？为此城作这样准备的人，不能回到他自己的城。啊，愚昧、反复无常且可怜的人！你们难道不明白这一切都属于他人，且在他人权下吗？因为这城的主人会说：\InnerQuote{我不愿意你住在我城里；离开此城，因为你不遵守我的法律。}你，虽然拥有田地、房屋和许多其他东西，但被他逐出时，你将怎样处理你的土地、房屋以及你为自己积蓄的其他财物？因为此地的主人公正地对你说：\InnerQuote{要么遵守我的法律，要么离开我的管辖。}那么，你打算怎样做？为了你的土地和其余财物，你拥有自己城的法律；你将完全否认你的法律，并按此城的法律行走。当心，否认你的法律会对你有害；因为如果你想回你的城，你将不被接纳，因为你否认了你城的法律，反被拒之门外。因此，要小心：如同一个居住在异乡的人，除仅够用的准备外，不要再为自己作更多准备；并且要准备好，当此城的主人因你违背他的法律而来驱逐你时，离开他的城，前往你自己的城，遵守你自己的法律，不受烦恼，而在大喜乐中。因此，你们这些事奉主、心中有主的人要当心，你们要作天主的作为，记念祂的诫命和祂所应许的许诺，并相信如果遵守祂的诫命，祂就必成就这一切。所以，不要购置田地，而要按各人的能力，购买受苦的灵魂，并探望寡妇和孤儿，不要忽视他们；并将你的财富和你从主所领受的一切准备，用于这些田地房屋。因为主人使你富足，正是为要你向祂行这些事；而购买这些田地、产业和房屋，远为更好，因为当你去居住时，你会在你自己的城里找到它们。这是一种高尚而神圣的花费，既无悲伤也无恐惧，却充满喜乐。不要效法外教人的花费，因为这对你们这些天主的仆人是有害的；却要实行你自己的花费，使你能在其中喜乐；也不要腐败，也不要触摸别人的东西，也不要贪恋，因为贪恋他人的财物是邪恶的；但要作你自己的工，你必得救。}\par

% structure: p0005 source=h2 target=subsection reason=relative-heading-rank; base=h2

\subsection{比喻二}

\Emph{正如葡萄树由榆树支撑，富人也借穷人的祈祷得到帮助。}\par

当我正在田间行走，观察一棵榆树和一棵葡萄树，并暗自思忖它们和它们的果实，牧者向我显现，对我说：\Quote{你正在想关于榆树和葡萄树的什么事呢？}\Quote{我正在想，}我回答说，\Quote{它们彼此非常相配。}\Quote{这两棵树，}他继续说，\Quote{是为天主的仆人树立的榜样。}\Quote{我想知道，}我说，\Quote{你说这两棵树所意图教导的榜样。}\Quote{你看见榆树和葡萄树了吗？}他说，\Quote{我看见了，先生，}我回答说。\Quote{这棵葡萄树，}他继续说，\Quote{结出果实，而榆树是不结果实的树；但是，除非葡萄树被缠绕在榆树上，它就不能在长时间延伸在地面上时结出许多果实；它所结的果实也是腐烂的，因为植株没有悬挂在榆树上。因此，当葡萄树被抛在榆树上时，它从自身和榆树都结出果实。而且你看，榆树也结出许多果实，不少于葡萄树，甚至更多；因为，}他继续说，\Quote{葡萄树当悬挂在榆树上时，结出许多好果实；但若抛在地上，它所结的果实又小又烂。所以，这个比喻是为天主的仆人——为穷人和富人设立的。}\Quote{先生，这是怎么回事？}我说；\Quote{请给我解释一下。}\Quote{听着，}他说：\Quote{富人拥有许多财富，但在关乎主的事上是贫穷的，因为他被自己的财富分心；他向主献上很少的宣认和代祷，他所献上的这些也是微小而软弱的，没有天上的力量。但是，当富人周济穷人，在他急需中帮助他，相信他对穷人所做的能在天主那里找到赏报——因为穷人富于代祷和宣认，他的代祷在天主面前大有能力——那时，富人毫不犹豫地帮助穷人一切事；穷人得到富人的帮助，就为他代祷，感谢天主恩赐与他的人。他还继续热切地关心穷人，使他的需要不断得到供应。因为他知道穷人的代祷在天主面前是蒙悦纳且有影响力的。于是，两者都完成他们的工作。穷人进行代祷；这是一项他所富有的工作，是他从主那里领受的，并用此回报帮助他的主人。同样，富人毫不迟疑地将他从主领受的财富施与穷人。这是一项伟大的工作，在天主面前蒙悦纳，因为他明白他财富的目的，将主的恩赐施与穷人，并正确地完成了对主的服务。然而，在人看来，榆树似乎不结果实，他们不知道也不明白：若遇干旱，含有水分的榆树滋养葡萄树；而葡萄树有了源源不断的水分，就结出双倍的果实，既为自己也为榆树。同样，穷人代祷主帮助富人，增加他们的财富；富人又在穷人急需中帮助他们，满足他们的灵魂。因此，两者在正义的工作中都是伙伴。做这些事的人必不被天主遗弃，而必被记在生命册上。拥有财富并明白它们是来自主的人是有福的。［因为凡有此心意的人，将能行一些善事。］}\par

% structure: p0008 source=h2 target=subsection reason=relative-heading-rank; base=h2

\subsection{比喻三}

\Emph{正如冬天绿叶树与枯树无法区分，同样，在此世义人与不义之人也无法区分。}\par

他向我显示许多没有叶子的树，看起来枯萎，如同我所见的；因为全都相似。他对我说：\Quote{你看见那些树了吗？}我回答说：\Quote{我看见了，先生，它们全都相似且枯萎。}他回答我说：\Quote{你所看见的这些树，是居住在这世上的人。}我说：\Quote{那么，先生，为什么它们看起来枯萎且相似呢？}他说：\Quote{因为，在此生，义人并不显明，罪人也不显明，而是相似的；因为今生对义人是冬天，他们不显现自己，因为他们与罪人同住；因为正如冬天里树木落叶后都相似，看不出哪些是死树哪些是活树，同样在这世上，义人也不显现，罪人也不显现，而是彼此相似。}\par

% structure: p0011 source=h2 target=subsection reason=relative-heading-rank; base=h2

\subsection{比喻四}

\Emph{正如夏天活树与枯树通过果实和活叶得以区分，同样，在来世，义人与不义之人在幸福中不同。}\par

祂又指给我看许多树木，有些正在发芽，有些已经枯干。祂对我说：\Quote{你看见这些树木吗？}我回答说：\Quote{我看见，先生，有些正在发芽，有些已经枯干。}祂说：\Quote{那些正在发芽的，是将在未来世界中生活的义人；因为未来的世界是义人的夏天，却是罪人的冬天。因此，当主的仁慈照耀时，那些天主的仆人必被显明出来，所有的人也必被显明出来。因为在夏天，每棵树的果子都显现出来，就能辨别它们是哪一种类；同样，义人的果子也必显现，所有在那个世界中结出果实的人都要被显明。至于外邦人和罪人，如同你所看见的枯干的树木，他们将是在那个世界中枯萎而不结果实的人，他们必如木柴被焚烧，并[因此]被显明出来，因为他们在生时行为邪恶。因为罪人必被消灭，因为他们犯了罪而没有悔改；外邦人必被焚烧，因为他们不认识创造他们的主。所以，你要结出果实，以便在那个夏天你的果实能被显明。并要避免过多的俗务，你就永不犯罪：因为那些忙于过多俗务的人也犯下许多罪，因他们的心被事务分心，丝毫不能事奉他们的主。}祂继续说：\Quote{那么，这样的人若不事奉祂，又如何能向主祈求并获得呢？事奉祂的人必获得他们的所求，但不事奉祂的人必一无所获。即便在履行一件小事上，人也能事奉主；因为他的心不会偏离主，而是以纯洁的心事奉祂。所以，你若行这些事，就能为未来的生命结出果实。而且，凡行这些事的人都必结出果实。}\par

% structure: p0014 source=h2 target=subsection reason=relative-heading-rank; base=h2

\subsection{比喻第五}

\Emph{论真守斋及其赏报，亦论身体的洁净。}\par

% structure: p0016 source=h3 target=subsubsection reason=relative-heading-rank; base=h2

\subsubsection{第一章}

我正守斋，坐在一座山上，为天主赐予我的一切恩惠向主谢恩时，看见牧者坐在我身旁，对我说：\Quote{你为什么一清早就到这里来了？}我回答说：\Quote{先生，因为我有一个\OriginalTerm{守斋}{station}.}祂问：\Quote{什么是守斋？}我回答：\Quote{先生，我正在守斋。}祂继续说：\Quote{你所守的这个斋是什么斋？}我回答说：\Quote{先生，我是按习惯守斋。}祂说：\Quote{你不知道如何向主守斋：你所守的这个无用的斋，对祂毫无价值。}我回答说：\Quote{先生，你为什么说这话？}祂继续说：\Quote{我对你说，你以为你守的那个斋不是斋。但我要教你什么是完全而中悦于主的斋。你听，}祂说：\Quote{天主并不想要这种空洞的守斋。因为用这种方式向天主守斋，你不能为正义的生活做什么；但你要向天主献上如下一种斋：在你的生活中不做邪恶的事，以纯洁的心事奉主；遵守祂的诫命，行在祂的律例中，不要让任何邪恶的欲望在你的心中升起；并要信靠天主。你若行这些事，敬畏祂，戒绝一切邪恶的事，你就要为天主而活；你若行这些事，你就守了一个伟大的斋，一个在天主面前可悦纳的斋。}\par

% structure: p0018 source=h3 target=subsubsection reason=relative-heading-rank; base=h2

\subsubsection{第二章}

\Quote{听我即将向你们讲述的关于守斋的比喻。\newline{}有一个人有一块田和许多奴隶，他在田里种植了一部分葡萄园，并挑选了一个忠实、可爱且极有价值的奴隶，叫他来，说：\InnerQuote{拿着这个我所种植的葡萄园，用桩子支起来直到我来，不要对葡萄园做别的事；要遵守我的命令，你就会从我这里得到自由。}\newline{}然后主人出国去了。\newline{}他走后，奴隶拿起并支起了葡萄园；当他把葡萄藤支好之后，他看到葡萄园里长满了杂草。\newline{}他思考说：\InnerQuote{我遵守了主人的命令；我要把葡萄园其余部分挖一下，挖过之后会更漂亮；并且清除杂草后，它会长出更多果实，不会被杂草窒息。}\newline{}于是他拿起工具挖了葡萄园，拔掉了里面所有的杂草。\newline{}那个葡萄园变得非常漂亮且多产，没有杂草窒息它。\newline{}过了一段时间，奴隶和那块地的主人回来了，进入葡萄园。\newline{}看到葡萄藤被桩子支持得很好，而且地面也被挖过，所有杂草被拔掉，葡萄藤多产，他对奴隶的工作非常满意。\newline{}他叫来他亲爱的儿子，即他的继承人，以及他的朋友们即他的谋士，告诉他们他给奴隶下了什么命令，以及他所发现的工作成果。\newline{}他们都为奴隶因主人所作的证言而欢喜。\newline{}他对他们说：\InnerQuote{我曾允诺这个奴隶，如果他遵守我给他的命令，就给他自由；他遵守了我的命令，并对葡萄园做了额外的好事，使我非常满意。因此，作为他所做工作的回报，我愿让他与我的儿子同为继承人，因为他有好的想法，没有忽视它们，而是实现了它们。}\newline{}主人的儿子和朋友们对这个决议很满意，即奴隶应与儿子同为继承人。\newline{}几天后，主人设宴，从他桌上送了许多菜肴给这个奴隶。\newline{}奴隶收到主人送来的菜肴后，取了自己足够的份，把剩下的分给其他奴隶。\newline{}其他奴隶很高兴收到菜肴，开始为他祈祷，希望他因这样对待他们而得到主人更大的恩宠。\newline{}主人听到了所发生的一切，再次对他的行为非常满意。\newline{}主人再次召集他的朋友和儿子，报告了奴隶关于他送去的菜肴的做法。\newline{}他们更加满意奴隶成为他儿子的共同继承人。}\par

% structure: p0020 source=h3 target=subsubsection reason=relative-heading-rank; base=h2

\subsubsection{第三章}

我对他说：\Quote{先生，我不理解这些比喻的含意，除非你向我解释，否则我无法明白。}\Quote{我将向你解释一切，} 他说，\Quote{并且凡我在我们交谈中所提到的，我都会向你展示。[遵守主的诫命，你将蒙悅纳，并被记在遵守他命令的人当中。]而且，如果你做了天主所命令之外的任何善事，你将为自己获得更丰盛的荣耀，并比否则更受天主的尊崇。因此，如果你在遵守天主的诫命之外，还做了这些服务，你若照我的命令遵守它们，便会有喜乐。}\newline{}我对他说：\Quote{先生，凡你吩咐我的，我必遵守，因为我知道你与我同在。}\newline{}\Quote{我必与你同在，} 他回答说，\Quote{因为你如此渴望行善；} 他又说，\Quote{我也要与所有有这样渴望的人同在。这个守斋，} 他接着说，\Quote{非常好，只要遵守了主的诫命。因此，你要这样遵守你所打算进行的守斋。首先，要提防一切恶言和一切恶欲，并从这世界的一切虚空中洁净你的心。如果你防备了这些事，你的守斋就是完美的。你还要这样做：在履行了所记载的事之后，在守斋的那一天，你只尝饼和水；并且算好那一天你本来打算吃的菜肴的价钱，把它送给一个寡妇、孤儿或某个有需要的人，这样你就表现出谦卑的心，使那从你的谦卑得益的人灵魂得饱足，并为主为你祈祷。如果你照我所命令的遵守守斋，你的祭献将为天主所悦纳，这守斋将被记录下来；如此所做的服务是尊贵的、神圣的，且为主所悦纳。因此，你要这样遵守，连同你的儿女和你全家，遵守这些事，你将有福；凡听见这些话并遵守的，也都有福；他们向主所求的，必得着。}\par

% structure: p0022 source=h3 target=subsubsection reason=relative-heading-rank; base=h2

\subsubsection{第四章}

我恳切地求他向我解释那田地的比喻、葡萄园的主人、那栽种葡萄园的仆人、木桩、从葡萄园拔出的杂草、儿子，以及作参议的朋友们，因为我知道这一切都是某种比喻。他回答我说：\Quote{你的问题真是固执。你本不应，}他继续说，\Quote{问任何问题；因为如果需要解释什么，它会让你知道的。}我对他说：\Quote{先生，无论你向我显示什么而不解释，我看见了也无益，因为我不明白它的意思。同样，如果你向我讲比喻而不阐释，我听了你的话也是枉然。}他再次回答我说：\Quote{每一个作为天主的仆人、心中怀着主的人，都向主祈求领悟，且必得着，并开启每一个比喻；凡以比喻说出的话，都为他所明白。但那些在祈祷中软弱懒惰的人，犹豫着不敢向主求什么；然而主满有怜悯，凡求他的，他都必赐予，从不落空。但你既被圣天使所坚固，从他得了这样的代祷，又不懒惰，为何不向主祈求领悟，从他领受呢？}我对他说：\Quote{先生，有你在我身边，我必须向你求问，因为你向我显示一切，并与我交谈；但若没有你，我看见或听见这些事，我便会求主解释它们。}\par

% structure: p0024 source=h3 target=subsubsection reason=relative-heading-rank; base=h2

\subsubsection{第五章}

\Quote{我刚才对你说过，}他回答说，\Quote{你在请求比喻的解释上狡猾而固执；但因你如此坚持，我将向你揭示那田地的比喻以及随后一切比喻的寓意，使你得以传扬给众人。现在倾听，}他说，\Quote{并明白它们。田地就是这个世界；田地的主就是那位创造、完善并坚固万物的；[子就是圣神；]仆人就是天主子；葡萄树就是他自己所栽种的这个民族；木桩就是主的圣天使，他们保守他的子民团结一起；从葡萄园拔出的杂草就是天主仆人的不义；他从筵席上派人送来的菜肴就是他通过自己的子赐给他子民的诫命；朋友和参议就是起初被造的圣天使；主人离家就是直到他显现之前所剩的时间。}我对他说：\Quote{先生，这一切都是伟大、奇妙而荣耀的事。我岂能，}我继续说，\Quote{明白它们？不能，任何其他人，即使极其聪明，也不能。}我又补充说：\Quote{请向我解释我正要问你的事。}\Quote{但说无妨，}他回答。\Quote{先生，}我问，\Quote{为何在比喻中天主子是仆人的形像？}\par

% structure: p0026 source=h3 target=subsubsection reason=relative-heading-rank; base=h2

\subsubsection{第六章}

\Quote{听，}他回答：\Quote{天主子不是仆人的形像，而是有大权能和大力。}\Quote{何以如此，先生？}我说；\Quote{我不明白。}\Quote{因为，}他回答，\Quote{天主栽种了葡萄园，也就是说，他创造了子民，并将他们交给他的子；子派他的天使看管他们；他自己涤净了他们的罪，经受了许多考验和劳苦，因为没有人能在没有劳动和辛劳的情况下挖掘。然后，他自己涤净了子民的罪后，将从他父领受的律法赐给他们，指示他们生命的道路。[你看，}他说，\Quote{他就是子民的主，从他父领受了全部权柄。] 并且主为何以他的子和荣耀天使为参议，关于仆人的继承权，请听。那神圣的、先在的圣神，创造了万有，天主使他居住在所拣选的肉身内。因此，这圣神所居住的肉身，高贵的服从于那圣神，敬虔而纯洁地行事，丝毫不玷污圣神；因此，在它卓越而纯洁地生活、与圣神一同劳苦合作、并在一切事上刚强勇敢地与圣神行动之后，他接纳了它作为自己的伙伴。因为肉身的这种行为使他喜悦，因为它在拥有圣神时没有在地上被玷污。因此，他以他的子和荣耀天使为同参议，为要使这毫无过犯地服从于身体的肉身拥有一个居所，并且不致使（它的服役的）报酬似乎丧失，因为那被发现毫无玷污、圣神居住其内的肉身，将获得报酬。你现在也得了这个比喻的解释。}\par

% structure: p0028 source=h3 target=subsubsection reason=relative-heading-rank; base=h2

\subsubsection{第七章}

\Quote{先生，我为听到这个解释而喜乐，}我说，\Quote{高兴。}\Quote{听，}他再次回答说：\Quote{你要保持这肉身纯洁无瑕，好使住在其中的圣神为你作证，你的肉身就得成义。要小心，绝不可让你的心里生出这样的念头：你的这肉身是可朽坏的，你就以任何亵渎的行为滥用它。如果你玷污了你的肉身，你也就玷污了圣神；如果你玷污了肉身[和灵魂]，你就不得存活。}\Quote{先生，如果有人，}我说，\Quote{在听到这些话之前一直无知，那玷污了自己肉身的人如何得救？}\Quote{至于以往无知所犯的罪，}他说，\Quote{唯有天主才能医治，因为一切权能都属于他。[但现在你要谨慎，全能而慈悲的天主将医治你以往的过犯，]如果你今后不再玷污你的身体和你的灵魂；因为二者是共同的，不能一个受玷污而另一个不受玷污：所以要保持二者纯洁，你就能为天主而生活。}\par

% structure: p0030 source=h2 target=subsection reason=relative-heading-rank; base=h2

\subsection{第六比喻}

\Emph{论两类享乐之人，以及他们的死、堕落，及其惩罚的期限。}\par

% structure: p0032 source=h3 target=subsubsection reason=relative-heading-rank; base=h2

\subsubsection{第一章}

我坐在家中，为所见到的一切光荣主，并思虑这些诫命：它们是卓越的、有力的、光荣的，并能拯救人的灵魂。我心中说：\Quote{如果我遵行这些诫命，我必蒙福；凡遵行它们的人也必蒙福。}当我正自言自语时，忽然看见他坐在我旁边，并听见他这样说：\Quote{你为什么对我所赐给你的诫命疑惑呢？它们是卓越的：你完全不要怀疑它们，但要怀有对主的信德，你就要遵行它们，因为我必在它们中坚固你。这些诫命有益于那些决意悔改的人：因为如果他们不遵行它们，他们的悔改便是枉然的。因此，你这悔改的人，要抛弃这世界耗尽你的邪恶；藉着穿戴圣善生活的一切美德，你将能遵守这些诫命，并且不再增添你的罪过。所以，你要遵行我的这些诫命，你就要为天主而活。这一切都是我向你说的。}他说完这些话后，又对我说：\Quote{让我们到田野去，我要给你看牧放羊群的牧人。}我回答说：\Quote{先生，我们去吧。}我们来到一块平原，他给我看一个年轻的牧人，身穿一套黄色衣服；他正在牧放许多羊，这些羊似乎在肥美地、放荡地、欢快地四处跳跃。那牧人自己因他的羊群而欢乐；牧人的神情欢快，他在羊群中跑来跑去。[我还看到其他一些羊在同一个地方放荡享乐，但却没有跳跃。]\par

% structure: p0034 source=h3 target=subsubsection reason=relative-heading-rank; base=h2

\subsubsection{第二章}

他对我说：\Quote{你看见这牧人吗？}我说：\Quote{我看见他，先生。}他回答说：\Quote{这是逸乐与欺骗的天使：他耗尽天主仆人们的灵魂，使他们偏离真理，以邪恶的欲望欺骗他们，他们因此丧亡；因为他们忘记生活天主的诫命，行走在欺骗和虚浮的逸乐中；他们被这天使毁灭，有的被带至死亡，有的被带至腐朽。}我对他说：\Quote{先生，我不明白这些话的意义——\InnerQuote{至死亡}和\InnerQuote{至腐朽}。}他说：\Quote{你听着。你所看见那欢快跳跃的羊，是那些永远离开天主、将自己交付给这世界的逸乐和欺骗[。在他们中间，没有藉着悔改回归生命的可能，因为他们除其他罪过外，还亵渎了主的名字。因此，这样的人被注定死亡。而你所看见那不跳跃、却在同一处吃草的羊，是那些将自己交付给逸乐和欺骗]，但并没有亵渎主的人。在他们中间还有悔改的希望，藉此得以生存。因此，腐朽尚有一种更新的希望，但死亡却是永远的毁灭。}\par

% structure: p0036 source=h3 target=subsubsection reason=relative-heading-rank; base=h2

\subsubsection{第三章}

我看见它们如此被打，被虐待，便为它们伤心，因为它们如此受折磨，完全没有休息。我对与我谈话的牧人说：\Quote{先生，这牧人是谁？如此冷酷严苛，完全不顾这些羊的怜悯？}他回答说：\Quote{这是惩罚的天使；他属于正义的天使，被任命施行惩罚。他因此带走那些偏离天主、行走在世界欲望和欺骗中的人，并按他们应得的，用可怕且多样的惩罚来责罚他们。}我说：\Quote{先生，我想知道这些多样刑罚和惩罚的性质是什么？}他说：\Quote{你听着，各种折磨和惩罚。这些折磨是在生活中发生的。因为有些人受损失惩罚，有些人受匮乏惩罚，有些人受各种疾病惩罚，还有些人受各种混乱与迷惑惩罚；另有人被不配的人侮辱，在许多其他方式下受苦：因为许多人，计划不稳定，尝试许多事，没有一样成功，他们说自己诸事不顺；但他们心中没有意识到自己做了恶事，反而责怪主。因此，当他们受尽各种苦难后，就被交付给我，为了良好的训练，他们在对主的信德中变得坚强；在余生中，他们以纯洁的心顺服主，诸事顺利，从主那里得到他们所求的一切；然后他们光荣主，因为他们被交付给了我，不再遭受任何灾祸。}\par

% structure: p0038 source=h3 target=subsubsection reason=relative-heading-rank; base=h2

\subsubsection{第四章}

我向他说：\Quote{先生，也请给我解释这个。}\Quote{你问什么？}他说。\Quote{先生，}我继续说，\Quote{那些沉迷于奢靡和被欺骗的人，是否被折磨的时间与他们沉迷奢靡和欺骗的时间相同？}他对我说：\Quote{他们受到同样的折磨。}[\Quote{先生，他们受的折磨要少得多，}我回答说；]\Quote{因为那些如此奢靡并忘记天主的人，应当受到七倍的折磨。}他对我说：\Quote{你糊涂了，不明白折磨的力量。}\Quote{为什么，先生，}我说，\Quote{如果我已经明白，就不会请你指示我了。}\Quote{听吧，}他说，\Quote{两者的力量。奢靡和欺骗的时间是一个时辰；但折磨的时辰相当于三十天。因此，如果一个人奢靡一天，被欺骗并受折磨一天，他受折磨的一天相当于一整年。所以，所有奢靡的日子，就有多少年的折磨要承受。由此可见，}他继续说，\Quote{奢靡和欺骗的时间非常短，但惩罚和折磨的时间却很长。}\par

% structure: p0040 source=h3 target=subsubsection reason=relative-heading-rank; base=h2

\subsubsection{第五章}

\Quote{可是，}我说，\Quote{关于欺骗、奢靡和折磨的时间，我仍然不完全明白；请给我更清楚地解释一下。}他回答说，并对我说：\Quote{你的愚昧是持久的；你也不愿净化你的心，并事奉天主。你要当心，}他补充说，\Quote{恐怕时间到了，你会被发现是愚昧的。现在听吧，}他补充说，\Quote{既然你渴望，你就能明白这些事。一个人奢靡一天，被欺骗一天，做他所愿做的事，就被许多愚昧所包裹，不明白他所做的事，直到第二天；因为他忘记了前一天所做的事。因为奢靡和欺骗没有记忆，因它们被愚昧所包裹；但当惩罚和折磨附着于人一天时，他受惩罚和折磨持续一年；因为惩罚和折磨有强大的记忆。因此，当他受折磨和惩罚一整年时，他终于想起他的奢靡和欺骗，并知道因这些他遭受苦难。因此，每一个奢靡和受欺骗的人都这样受折磨，因为虽然他们有生命，却把自己交给了死亡。}\Quote{先生，哪些奢靡是有害的？}我问。\Quote{一个人的每一个以快乐完成的行为，}他回答说，\Quote{都是奢靡行为；因为急躁的人，当他满足自己的倾向时，就在奢靡；奸淫者、酗酒者、毁谤者、说谎者、贪婪者、偷窃者以及做类似这些事的人，都满足了自己特有的癖好，并因此奢靡。所有这些奢靡行为对天主的仆人们都是有害的。因此，因这些欺骗，那些受惩罚和折磨的人就受苦。也有一些奢靡能拯救人；因为许多行善的人沉迷于奢靡，被自己的快乐所吸引：然而这种奢靡对天主的仆人们是有益的，并为这样的人赢得生命；但前面列举的有害奢靡行为给他们带来折磨和惩罚；如果他们继续这些行为而不悔改，就给自己带来死亡。}\par

% structure: p0042 source=h2 target=subsection reason=relative-heading-rank; base=h2

\subsection{第七比喻}

\Emph{悔改的人必须结出与悔改相称的果实。}\par

过了几天，我在同一片平原上看见了他——就是从前我看见那些牧人的地方。他对我说：\Quote{你对我有什么要求？}我对他说：\Quote{先生，请你命令那位惩罚的牧人离开我的家，因为他极其折磨我。}他回答说：\Quote{你必须受折磨；因为，}他接着说，\Quote{光荣的天使曾这样命令关于你的事，他愿意你被考验。}我回答说：\Quote{先生，我做了什么如此邪恶的事，竟要被交给这位天使？}他说：\Quote{听着：你的罪很多，但还不至于严重到需要把\Emph{你}交给这位天使；但你的家眷犯了极大的不义和罪恶，光荣的天使因他们的行为而向他们发怒；因此他命令你受一段时期的折磨，好使他们也能悔改，并净化自己脱离这世界的一切欲望。因此，当他们悔改并得到净化时，那位惩罚的天使就会离去。}我对他说：\Quote{先生，如果他们做了那样的事以致光荣的天使向他们发怒，那么我又做了什么？}他回答说：\Quote{他们根本无法受折磨，除非你——家长——受折磨；因为你受折磨时，他们必然也受折磨；但如果你安舒，他们就感觉不到折磨。}我说：\Quote{先生，他们已全心全意悔改了。}他回答说：\Quote{我也知道他们已全心全意悔改了；但是你以为悔改者的罪就被赦免了吗？并非完全如此，但悔改者必须折磨自己的灵魂，在一切行为上极其谦卑，受多种折磨；如果他忍受临到他身上的折磨，那创造万有并赋予它们能力的主必定会怜悯他，并医治他；祂看见每个忏悔者的心纯洁远离一切邪恶时，就会这样做；而现在你和你的家受折磨是有益的。但我何必对你说许多呢？你必须受折磨，如同那位把你交给我的主的天使所命令的。为此，你要感谢主，因为祂认为你配得预先知道这折磨，好使你在它来临之前就知道，从而能勇敢地承受它。}我对他说：\Quote{先生，请与我同在，我便能承受一切折磨。}他说：\Quote{我必与你同在，并且我会请求惩罚的天使更轻地折磨你；尽管如此，你仍要受一段短暂的折磨，然后你将被重新安置在你的家中。只要继续谦卑，并全心全意纯洁地事奉主，你和你的儿女，以及你的家，并遵行我所命令你的诫命，你的悔改就会深刻而纯洁；如果你和你的家人遵守这些事，一切折磨就会离开你。}他又说：\Quote{折磨也必离开所有遵行我这些诫命的人。}\par

% structure: p0045 source=h2 target=subsection reason=relative-heading-rank; base=h2

\subsection{第八个比喻}

\Emph{被选者与忏悔者的罪有多种，但所有人将按他们悔改和善行的程度获得酬报。}\par

% structure: p0047 source=h3 target=subsubsection reason=relative-heading-rank; base=h2

\subsubsection{第一章}

他给我看一棵巨大的柳树，遮蔽了平原和山岳，所有奉主名被召的人都聚集在这棵柳树的荫下。一位光荣的主的天使，身材极高，站在柳树旁，手中拿着一把大修枝刀，他从柳树上砍下小枝，分给那些被柳树荫蔽的人；他给他们的枝子很小，大约一肘长。他们全都接过枝子后，天使放下修枝刀，那棵树依然完好，如同我起初所见。我心中诧异，说道：\Quote{既然砍下这么多枝条，树怎会完好呢？}牧人对我说：\Quote{不要惊讶树在砍下许多枝条后仍然完好；［但请等待，］等你看到一切，就会向你解释其含义。}那位将枝条分给众人的天使又向他们索要枝条，按他们接受枝条的顺序被召到他面前，每人归还自己的枝条。主的天使接过并察看它们。从有些人那里，他收到枯萎且被虫蛀的枝条；那些归还这种枝条的人，主的天使命令他们站在一边。另一些人归还的枝条枯萎但未被虫蛀；他命令这些人站在一边。又有些人归还半枯萎的枝条，这些人站在一边；又有些人归还半枯萎且有裂纹的枝条，这些人站在一边。［又有些人归还绿色且有裂纹的枝条，这些人站在一边。］又有些人归还的枝条一半枯萎一半绿色，这些人站在一边。又有些人带来枝条三分之二绿色，三分之一枯萎，这些人站在一边。又有些人归还的枝条三分之二枯萎，三分之一绿色，这些人站在一边。又有些人归还的枝条几乎全绿，只有极小部分——顶端——枯萎，但有裂纹，这些人站在一边。另有一些人只有极少部分绿色，其余部分枯萎，这些人站在一边。又有些人带来枝条，如同从天使手中接到时那样青绿。大多数群众归还这类枝条，天使对此极为喜悦，这些人站在一边。［又有些人归还青绿且有嫩芽的枝条，这些人站在一边，天使因他们而极为喜悦。］又有些人归还青绿且有嫩芽的枝条，嫩芽上仿佛有些果实；那些枝条被发现是这种类型的人极其欢喜。天使因他们而欢欣；牧人也为他们大大喜乐。\par

% structure: p0049 source=h3 target=subsubsection reason=relative-heading-rank; base=h2

\subsubsection{第二章}

主的天使命令取来冠冕；冠冕取来了，仿佛是用棕榈枝编成的；他给那些交回带有嫩枝和果实的树枝的人戴上冠冕，并打发他们进入塔中。其余的人，即那些交回绿色、有嫩枝但无果实的树枝的人，他也赐给他们印鉴，打发他们进入塔中。所有进入塔中的人都穿着同样的衣服——洁白如雪。至于那些交回的树枝仍是绿的、如同当初领受时一样的人，他释放了他们，赐给他们衣服和印鉴。天使做完这些事之后，对牧者说：\Quote{我要走了，你要按照各人配得住所的程度，把这些人都送进墙内。你要仔细检查他们的树枝，然后才打发他们走；但是要\Emph{仔细}检查。}他接着说：\Quote{小心不要让他们有人逃脱你；如果有人逃脱你，我将在祭台前审问他。}对牧者说完这些话，天使就离开了。天使离开后，牧者对我说：\Quote{我们拿所有这些人的树枝去栽种，看看其中有没有能活的。}我对他说：\Quote{先生，这些枯枝怎么能活呢？}他回答说：\Quote{这棵树是柳树，是一种生命力很强的树种。所以，如果这些树枝被栽种，得到一点水分，许多就会活过来。现在我们试试，把水浇在上面；如果其中有活的，我就与他们一同喜乐；即使不活，至少我不会被认为疏忽。}牧者又命我按各人所站的位置召唤他们。他们一排一排地前来，将树枝交给牧者。牧者接过树枝，一行一行地栽种；栽种完后，他浇了许多水在上头，以致水把树枝都淹没了；树枝吸收了水分之后，他对我说：\Quote{我们走吧，过几天再回来察看所有的树枝；因为创造这棵树的那一位愿意所有从它领受树枝的人都活。我也希望这些吸收了水分并喝了水的树枝，大部分都能活。}\par

% structure: p0051 source=h3 target=subsubsection reason=relative-heading-rank; base=h2

\subsubsection{第三章}

我对他说：\Quote{先生，请给我解释这棵树的意义，因为我对此感到困惑；那么多树枝被砍下之后，它仍然完好，看不出有被砍去的痕迹。这令我困惑。}他说：\Quote{你听着：这棵大树的影子遮蔽平原、山岭和整个大地，它就是天主颁给全世界的法律；这法律就是天主子，被传扬到地极；在它荫下的人民，就是听了这宣讲并相信他的人。那伟大而荣耀的天使\Person{弥额尔}是这人民的掌权者，治理他们；因为他把法律放在信者的心里；他监察那些领受法律的人，看他们是否遵守了。你看见各人的树枝，因为树枝就是法律。你看见许多无用的树枝，你将认出他们全部——那些不遵守法律的人；你也将看见各人的住所。}我对他说：\Quote{先生，为什么他打发一些人进入塔中，却把另一些人留给你呢？}他回答说：\Quote{所有违背了从他领受的法律的人，他都留在我权下，为使他们悔改；但所有满足了法律并遵守它的人，他留在他自己的权下。}我接着问：\Quote{那么，那些戴冠冕并进入塔中的人是谁呢？}\Quote{他们是因法律而受苦的人；至于其他人，那些交回绿色、有嫩枝但无果实的树枝的，是因法律而受磨难，但没有受苦，也没有否认法律的人；那些交回的树枝仍是绿的、如同当初领受时一样的人，是令人尊敬的人、义人，以及那些以纯洁之心谨慎行走、遵守了主的诫命的人。其余的人，等我察看了那些栽种并浇灌的树枝之后，你就会知道。}\par

% structure: p0053 source=h3 target=subsubsection reason=relative-heading-rank; base=h2

\subsubsection{第四章}

过了几天，我们来到那地方，牧者坐在天使的位置上，我站在他旁边。他对我说：\Quote{你要用纯净的、未漂白的粗麻布束腰；}见我束好了腰，准备服事他，他说：\Quote{你召唤那些被栽种的树枝所属的人，按照各人当初交树枝的顺序。}于是我走到平原，召唤他们全体，他们都按行列站好。他对他们说：\Quote{每人拔出自己的树枝，带到我这里来。}首先交出树枝的是那些树枝枯干且被砍断的人；因为发现它们如此枯干且被砍断，他命令他们站在一旁。接着那些人交出树枝，他们的树枝枯干但没有被砍断。有些人交回绿色的树枝，有些人交回枯干且貌似被虫蛀过的树枝。于是他命令交回绿色树枝的人站在一旁；交回枯干且被砍断的树枝的人，他命令他们与第一批人站在一起。接着交出树枝的是那些树枝半枯干且有裂口的人；他们中许多人交回绿色无裂口的树枝；也有一些是绿色、带嫩枝且嫩枝上有果实的，正如那些戴冠冕后进入塔中的人所交回的。有些人交出枯干且被虫蛀过的树枝，有些人交出枯干且未被虫蛀的；有些人的树枝半枯干且有裂口。他命令各人分开站立，一些人朝向自己的行列，另一些人则与他们分开。\par

% structure: p0055 source=h3 target=subsubsection reason=relative-heading-rank; base=h2

\subsubsection{第五章}

然后那些枝子仍青但已破裂的人交上了枝子：所有这些人都交上青枝，并站在自己的行列里。牧者对此感到满意，因为他们全都改变了，裂痕消失了。那些枝子半青半枯的人也交上了枝子：有些人的枝子完全青翠，有些人的半枯，有些人的枯干且被虫蛀，有些人的青翠且带有嫩枝。所有这些人都被送走，各自到自己的行列里。[接下来，那些枝子三分之二青、三分之一枯的人交上了枝子。他们中许多人的枝子半枯；另一些枯干且腐烂；另一些半枯且有裂痕，少数是青的。所有这些人都站在自己的行列里。]那些枝子青翠但只有极轻微枯干和裂痕的人也交上了枝子。其中有些人交上青枝，有些人交上青翠且带有嫩枝的枝子。这些人也去到自己的行列里。接着，那些只有很小部分青翠、其余部分枯干的人交上了枝子。这些人的枝子大部分被发现是青翠且带有嫩枝，嫩枝上还有果实，另一些则完全青翠。牧者对这些枝子极其满意，因为它们被发现处于这种状态。这些人也离去，各自到自己的行列里。\par

% structure: p0057 source=h3 target=subsubsection reason=relative-heading-rank; base=h2

\subsubsection{第六章}

牧者查考了所有人的枝子后，对我说：\Quote{我告诉过你，这树生命力坚韧。你看，}他接着说，\Quote{多少人悔改了并得救了。}\Quote{我看见，先生，}我答道。\Quote{为让你看见，}他补充说，\Quote{主的极大仁慈，是何等伟大而荣耀，祂已将祂的圣神赐予那些堪当悔改的人。}\Quote{那么，先生，}我说，\Quote{为什么不是所有这些人都悔改了呢？}他答道：\Quote{凡祂看透其心将变得纯洁顺服祂的人，祂赐予他们全心悔改的能力；但凡祂察觉其诡诈和邪恶、并看见他们打算假冒伪善地悔改的人，祂就不赐予悔改，免得他们再次亵渎祂的名。}我对他说：\Quote{先生，现在请向我指明那些交上枝子的人，他们是怎样的人，他们的住所如何，好让那些曾经相信、领受了印记、却又毁坏了它、未能完好保存的人，在认识到自己的行为后，能够悔改，并从你那里领受一个印记，并荣耀主，因为祂怜悯了他们，并派遣你来更新他们的心神。}\Quote{请听，}他说：\Quote{那些枝子被发现枯干且被虫蛀的人，是教会的叛教者和背信者，他们在自己的罪中亵渎了主，并且以他们被呼召的主的名号为耻。因此，他们最终失落于天主。你看，他们中没有一个人悔改，尽管他们听到了我对他们所说的话，就是我嘱咐你的那些话。生命已从这样的人身上离去。那些交上枯干而未朽坏枝子的人，他们也与前者相近；因为他们是伪善者，是引入异端学说的人，是颠覆天主仆人的人，尤其是颠覆那些犯了罪的人，不让他们悔改，反而用愚拙的学说劝诱他们。这些人在悔改上仍有希望。你看，自从我对他们说话以来，他们中许多人也已悔改，并且还会继续悔改。但所有不愿悔改的人都丧失了生命；而他们中悔改的许多人都变好了，他们的住所被指定在第一层墙内；有些人甚至登上了塔。那么你看，}他说，\Quote{悔改带给罪人生命，而不悔改则带来死亡。}\par

% structure: p0059 source=h3 target=subsubsection reason=relative-heading-rank; base=h2

\subsubsection{第七章}

\Quote{至于那些交上枝子半枯且有裂痕的人，也听关于他们的事。那些枝子半枯程度相同的人是动摇者；因为他们既不算活着，也不算死了。那些枝子半枯且有裂痕的人，既是动摇者又是诽谤者，[诋毁不在场的人，]彼此永不和睦，总是纷争不息。然而对这些也是如此，}他继续说，\Quote{悔改是可能的。你看，}他说，\Quote{他们中有些人已经悔改，而且他们里面仍然存留着，}他接着说，\Quote{悔改的希望。至于他们中，}他补充道，\Quote{凡已悔改的，将在塔内得住所；那些悔改较慢的将住在墙内；凡根本不悔改、仍行其事的，必要全然灭亡。那些交上青绿但有裂痕枝子的人，总是忠信而良善，却为着首位和名声彼此争竞：如今所有这些人都愚昧，沉溺于这样的竞争。然而他们，因本性良善，一听见我的诫命，就洁净了自己，很快悔改了。因此，他们的住所是在塔内。但若有人再陷入纷争，他必被逐出塔外，丧失生命。生命属于所有遵守主诫命的人；但在诫命中，没有关于首位或任何荣耀的竞争，只有关于忍耐和个人谦卑的竞争。那么，这样的人身上有主的生命，而在好争竞者和违背者身上，则是死亡。}\par

% structure: p0061 source=h3 target=subsubsection reason=relative-heading-rank; base=h2

\subsubsection{第八章}

\Quote{那些交出半绿半枯树枝的，是沉浸于世俗事务、不亲近圣徒的人。因此，他们一半是活的，一半是死的。许多听了我的命令的人悔改了，至少那些悔改的人住进了塔中。但其中有些人最终堕落了：这些人没有悔改，因为他们因自己的事务亵渎了主，并否认了祂。他们因此因所行的恶而丧失了生命。他们中许多人犹豫不决。这些人仍有权悔改，只要他们尽快悔改；他们的住所将在塔中。但如果他们悔改迟缓，他们将住在墙内；如果他们不悔改，他们也丧失了生命。那些交出三分之二枯、三分之一绿树枝的，是以各种方式否认主的人。然而，许多人悔改了，但其中有些人犹豫不决、心存疑惑。因此，这些人仍有机会悔改，只要他们尽快悔改，不耽于享乐；但如果他们执迷于自己的行为，这些人也给自己招致死亡。}\par

% structure: p0063 source=h3 target=subsubsection reason=relative-heading-rank; base=h2

\subsubsection{第九章}

\Quote{那些交出三分之二枯、三分之一绿树枝的，是那些曾经忠信的人；但在致富并在外邦人中成名后，他们披上极大的骄傲，变得心高气傲，离弃了真理，不再亲近义人，反而与外邦人同住，这种生活方式更令他们欢喜。他们没有离开天主，仍保持信德，却不实践信德的行为。他们中许多人因此悔改了，他们的住所在塔中。另一些人继续与外邦人生活直到终了，被虚浮的荣耀所败坏，离开了天主，事奉外邦人的工作和行为。这些人被算作外邦人。但还有一些人犹豫不决，因自己所行的事而不指望得救；另一些人心存疑惑，在自己中间引起分裂。因此，那些因自己的行为而疑惑的人，悔改之门仍然敞开；但他们的悔改应当迅速，这样他们的住所才能在塔中。至于那些不悔改、耽于享乐的人，死亡临近了。}\par

% structure: p0065 source=h3 target=subsubsection reason=relative-heading-rank; base=h2

\subsubsection{第十章}

\Quote{那些交出绿色但尖端枯萎开裂的树枝的，这些人向来善良、忠信、在天主面前尊贵；但他们因放纵小欲望而犯了小罪，并彼此挑剔。但听了我的话后，他们中大部分人迅速悔改了，他们的住所在塔上。然而有些人仍犹豫不决；其中一些疑惑的人造成了更大的分裂。因此，这些人仍有悔改的希望，因为他们向来善良；他们中几乎无人会灭亡。那些交出枯枝、只有极小部分绿色的，是那些只相信、却继续行不义之事的人。他们从未离开天主，而是欣然背负祂的名，并喜乐地接待祂的仆人进入自己的家。他们因此听了这悔改之道后，毫不犹豫地悔改了，并实践一切德行和义行；有些人甚至受苦，甘愿被处死，因为他们知晓自己所行的事。因此，所有这些人的住所将在塔中。}\par

% structure: p0067 source=h3 target=subsubsection reason=relative-heading-rank; base=h2

\subsubsection{第十一章}

他解释完所有树枝后，对我说：\Quote{你去告诉每一个人，叫他们悔改，他们就能为天主而活。因为主怜悯了众人，打发我来赐予悔改，尽管有些人因自己的行为不配；但主忍耐宽容，愿意那些被圣子召叫的人得救。} 我对他说：\Quote{先生，我希望所有听见这些话的人都会悔改；因为我确信，每个人一旦认识到自己的行为，并敬畏主，就会悔改。} 他回答我说：\Quote{凡全心从先前列举的邪恶中洁净自己，不再加添罪过的人，若对这些诫命不犹豫，必从主获得以往过犯的医治；他们将为天主而活。但你当遵行我的诫命，并活着。} 他指示了这些事，说了这一切话后，对我说：\Quote{其余的事，过几天我会指示你。}\par

% structure: p0069 source=h2 target=subsection reason=relative-heading-rank; base=h2

\subsection{第九比喻}

\Emph{战斗与得胜的教会建造中的伟大奥秘。}\par

% structure: p0071 source=h3 target=subsubsection reason=relative-heading-rank; base=h2

\subsubsection{第一章}

在我写完了牧者——补赎天使——的诫命和比喻之后，他来到我面前说：\Quote{我希望向你解释，那以教会的形像与你说话的圣神所指示你的事，因为那圣神就是天主子。因为你肉身有些软弱，天使没有向你解释。然而，当你被圣神坚固、力量增强，以至于你也能看见天使时，教会便将塔楼的建造指示给你。你以一种高贵而庄严的方式看见了这一切，如同由一位贞女所指示；但现在你藉着同一位圣神，如同由一位天使所指示而看见。然而，你必须从我这里更准确地学习一切。因为光荣的天使为此目的派遣我住在你家中，使你能够带着能力看见一切，毫无恐惧，就像从前一样。} 然后他带我前往\Place{阿卡迪亚}，到一个圆形的山丘；他把我放在山顶上，给我看一片广阔的平原，平原周围有十二座山，形状各不相同。第一座山黑如煤烟；第二座光秃秃的，没有草；第三座长满了荆棘和蒺藜；第四座草木半枯，上部的植物是绿的，根部的部分枯萎了；有些草被太阳晒焦后便枯萎了。第五座山有绿草，但崎岖不平。第六座山布满了裂缝，有些小，有些大；裂缝中有草，但植物不很茂盛，反倒像是朽坏了。第七座山有喜乐的牧场，整座山繁花盛开，各种牲畜和飞鸟都在那座山上吃草；牲畜和鸟吃得越多，那座山上的草就越茂盛。第八座山布满了泉源，主的各种受造物都饮用那座山的泉水。但第九座山[根本没有水，完全是沙漠，里面有致人死命的毒蛇。第十座山]有非常大的树，完全被荫蔽，在树荫下羊群躺着休息和反刍。第十一座山树木非常茂密，那些树结果实，装饰着各种水果，使任何看见的人都想吃它们的果实。而第十二座山则是全白的，面貌喜乐，山本身非常美丽。\par

% structure: p0073 source=h3 target=subsubsection reason=relative-heading-rank; base=h2

\subsubsection{第二章}

在平原中央，他给我看一块从平原升起的巨大白色岩石。这块岩石比那些山更高，呈矩形，能够容纳整个世界；那块岩石是古老的，其上开凿了一道门；门的开凿在我看来像是新近完成的。门在阳光下闪耀到如此程度，使我惊叹于门的光辉；门周围站着十二位贞女。站在角落的四位在我看来比其他人更出众——尽管她们全部都很出众——她们站在门的四边；每边之间有两位贞女。她们穿着亚麻长袍，优雅地束着腰带，露出右肩，仿佛要承担什么重担。她们就这样站着，准备就绪；因为她们极其喜乐和热切。看见这些后，我内心惊奇，因为我正看见伟大而光荣的景象。我又因那些贞女而困惑，她们虽然如此娇弱，却勇敢地站着，仿佛要承担整个上天。牧者对我说：\Quote{你为什么心里推理、困惑、折磨自己呢？因为你不能理解的事，不要试图去理解，仿佛你聪明似的；但要向主祈求，使你能获得领悟并认识它们。你看不见后面的事，却看见前面的事。因此，凡你看不见的，就放下它，不要为此折磨自己；但你所看见的，要掌握它，不要在其他事上浪费精力；我将向你解释我显示给你的一切。所以，看那些剩下的事。}\par

% structure: p0075 source=h3 target=subsubsection reason=relative-heading-rank; base=h2

\subsubsection{第三章}

我看见六个人走来，高大、出众、样貌相似，他们召唤了一大群人。来的人也是高大、俊美、强壮的人；那六个人命令他们在岩石上建造一座塔楼。来建造塔楼的人发出巨大的响声，他们绕着门跑来跑去。站在门周围的贞女们告诉那些人赶快建造塔楼。贞女们伸出双手，仿佛要从那些人那里接受什么。那六个人命令石头从一个坑中升上来，并去建造塔楼。有十块闪亮的长方形石头升上来，未经采石场打磨。那六个人叫来贞女，吩咐她们拿起所有用于建造的石头，穿过门，交给将要建造塔楼的人们。贞女们将最初从坑中升上来的十块石头互相叠放，一起搬运，每块石头单独搬运。\par

% structure: p0077 source=h3 target=subsubsection reason=relative-heading-rank; base=h2

\subsubsection{第四章}

当他们一起站在门周围时，那些看来强壮的人就扛起它们，并在石头角下弯腰；其他人在石头边下弯腰。就这样他们把所有石头都搬走了。然后他们按照命令经过门搬运，交给建筑塔的人；那些人接过石头继续建造。现在，塔建在大岩石上，并在门的上方。那十块石头被预备作塔的基础。岩石和门是整座塔的支撑。在那十块石头之后，又有其他二十[五]块从坑中出来，这些被童贞女像之前一样搬去，嵌入塔的建筑中。之后，三十五块升上来。这些也以同样方式嵌入塔中。之后，又有其他四十块石头上来；所有这些都被投入塔的建筑中，[在塔的基础中有四行，]他们停止从坑中上升。建筑者也暂停了一会儿。然后，那六个人命令人群中的大众从山上搬运石头用于塔的建筑。于是，从各座山运来了各种颜色的石头，经过那些人切割后交给童贞女；童贞女经过门运进去，交给建筑塔用。当各种颜色的石头被放入建筑时，它们都变得一样白，失去了各自颜色。有些石头被那些人交给建筑用，这些没有变得闪亮；但正如它们被放置时那样，它们被发现仍然保持原样：因为它们不是由童贞女给的，也不是经过门运进的。因此，这些石头在塔的建筑中与其他石头不协调。那六个人看到建筑中有这些不合适的石头，就命令把它们搬走，搬回它们被取来的地方；[一块一块地移开，放在一边；]他们对运石头的人说，\Quote{不要带任何石头来建筑，只要把它们放在塔旁，以便童贞女经过门运进去，交给建筑用。因为除非，}他们说，\Quote{它们经过门由童贞女的手运送，否则不能改变颜色：因此，不要徒劳，}他们说，\Quote{毫无用处。}\par

% structure: p0079 source=h3 target=subsubsection reason=relative-heading-rank; base=h2

\subsubsection{第五章}

那天建筑完成了，但塔没有完工；因为还要再添加建筑，建筑工作暂停了。那六个人命令所有建筑者后退一点休息，但吩咐童贞女不要离开塔；在我看来，童贞女被留下看守塔。现在，当所有人都退下休息时，我对牧者说，\Quote{为什么塔的建筑没有完成？}牧者回答，\Quote{塔，}他回答，\Quote{现在还不能完成，直到它的主来检查建筑，以便若发现任何石头变质，他可以把它们换掉；因为塔是按照他的意愿建造的。}\Quote{先生，我想知道，}我说，\Quote{这座塔的建筑是什么意思，以及岩石、门、山、童贞女、从坑中上来的未经切割就进入建筑的石头都代表什么。为什么起初在基础上放十块石头，然后二十五块，然后三十五块，然后四十块？我还想知道关于那些进入建筑又被取出送回原处的石头的事情。先生，请在这些问题上使我安心，并向我解释。}他回答说，\Quote{如果你被发现不是对琐事好奇，}他回答，\Quote{你就会知道一切。因为过几天[我们将来到这里，你将看到发生在这座塔上的其他事情，并准确知道所有比喻。}过几天]我们来到坐下的地方。他对我说，\Quote{我们去塔那里；塔的主人要来检查它。}我们来到塔那里，除了童贞女，附近没有别人。牧者问童贞女塔主人是否来了；他们回答说他就要求检查建筑。\par

% structure: p0081 source=h3 target=subsubsection reason=relative-heading-rank; base=h2

\subsubsection{第六章}

看哪，过了一会儿，我看见一大群人来了，中间有一个人身材如此高大，高过塔。曾参与建筑的六个人与他在一起，还有许多其他尊贵的人在他周围。看守塔的童贞女跑上前去亲吻他，开始在他旁边绕着塔走。那个人仔细检查建筑，逐一触摸每块石头；他手里拿着一根棍子，每块石头敲了三下。当他敲打时，有些石头变得像煤灰一样黑，有些看似布满痂疤，有些开裂，有些残缺，有些既不白也不黑，有些粗糙且与其他石头不协调，有些有[非常多的]污渍：这就是在建筑中发现的变质石头的多样性。他命令所有这些石头从塔中取出，放在塔旁，并带其他石头来替换。[建筑者问他希望从哪座山上取石头来替换。]他没有命令从山上取石头，[而是命令从附近一片平原上取。]平原被挖开，发现了闪亮的长方形石头，也有一些圆形的；平原上所有石头都被运来，由童贞女经过门搬运进去。长方形石头被切割后放在被取走石头的原处；但圆形的石头没有被放入建筑，因为它们难以切割，似乎不易接受凿子；但它们被存放在塔旁，似乎是要被切割后用于建筑，因为它们极其明亮。\par

% structure: p0083 source=h3 target=subsubsection reason=relative-heading-rank; base=h2

\subsubsection{第七章}

那位光荣的人，全塔之主，完成了这些改动之后，便召牧人前来，把所有堆在塔旁、从建筑中弃置的石头都交给他，对他说：\Quote{你要仔细清洁所有这些石头，把那些能与其余石头谐合以用于建塔的挑出来；不能谐合的，就远远地抛离塔旁。}[将这样的命令交给牧人之后，他就与所有同来的人一起离开了塔。]那时，童贞女们正站在塔四周看守着塔。我又对牧人说：\Quote{这些石头被弃置之后，还能再回到建塔工程中吗？}他回答我说：\Quote{你看见这些石头了吗？}\Quote{我看见了，先生，}我答道。\Quote{这些石头的大部分，}他说，\Quote{我要凿刻，放进建筑中，它们会与其余的谐合。}\Quote{先生，}我说，\Quote{它们被周身切割之后，怎么能填满同样的空间呢？}他答道：\Quote{那些发现是小的，将被抛进建筑的中心；那些较大的，将被放在外面，它们会把小石头夹住。}说完这些话，他对我说：\Quote{我们走吧，两天后我们再来清洁这些石头，并把它们投进建筑；因为塔周围的一切都必须清洁，免得主人突然到来，发现塔四周的地方肮脏而不悦，以致这些石头不能送回建塔，而我也显得对主人疏忽了。}两天后，我们来到塔前，他对我说：\Quote{我们查验所有的石头，确定哪些可以送回建筑。}我对他说：\Quote{先生，我们查验吧！}\par

% structure: p0085 source=h3 target=subsubsection reason=relative-heading-rank; base=h2

\subsubsection{第八章}

我们从黑石头开始首先查验。那些从建筑中取出的石头，发现仍然保持原状；牧人命令将它们从塔旁移开，另放一处。接着他查验那些有疤痂的石头；他拿起其中许多进行凿刻，命令童贞女们拿起并投进建筑。童贞女们举起它们，放在塔建筑的中心。其余的他命令放在黑石头旁边；因为发现这些也是黑的。接着他查验那些有裂缝的石头；他凿刻了许多，命令童贞女们搬运到建筑中：它们被放在外面，因为发现它们比别的更坚固；但其余的，因裂缝众多不能凿刻，为此它们便被弃于建塔工程之外。接着他查验那些有缺口的石头，发现其中许多是黑的，有些则有大的裂缝。这些他也命令与那些已被弃置的放在一起。但其余的，经过清洁和凿刻后，他命令放进建筑。童贞女们拿起它们，安放在塔建筑的中心，因为它们有点软弱。接着他查验那些半白半黑的石头，发现其中许多是黑的。他也命令将这些与那些已被弃置的一起拿走。其余的全部由童贞女们拿走；因为它们是白的，由童贞女们自己安放进建筑。它们被放在外面，因为发现它们坚固，能支撑中心放置的那些，因为它们完全没有缺口。接着他查验那些粗糙坚硬的石头；其中少数被弃置，因为无法凿刻，发现它们极其坚硬。但其余的经过凿刻，由童贞女们搬走，安放在塔建筑的中心；因为它们有点软弱。接着他查验那些有污渍的石头；其中极少数是黑的，被与其他的一起扔掉；但大部分被发现是明亮的，这些由童贞女们安放进建筑，但因它们的坚固被放在外面。\par

% structure: p0087 source=h3 target=subsubsection reason=relative-heading-rank; base=h2

\subsubsection{第九章}

他接着前来察看那些白色而圆润的石头，对我说：\Quote{这些石头我们该做什么？} \Quote{我怎么会知道，先生？}我答道。\Quote{你难道没有打算吗？} \Quote{先生，}我回答，\Quote{我不懂这门手艺，也不是石匠，无法说。} \Quote{莫非你没有看见，}他说，\Quote{这些石头极其圆润？若我要使它们成为矩形，必须切掉很大一部分；因为其中有些必须被放进建筑中。} \Quote{既然如此，}我说，\Quote{它们必须如此，你又何必自寻烦恼，不立即挑选你中意的石头来建造，并安置进去呢？}他便从中选出了较大的和发亮的石头，加以切割；贞女们则搬起它们，安置在建筑的外部。其余剩下的被搬走，放在它们被取来的平原上。然而，它们并非被抛弃，\Quote{因为，}他说，\Quote{还要为塔再加建一小部分。这塔的主人愿意所有石头都被安置在建筑中，因为它们极其明亮。}于是召来了十二位妇女，容貌极美，身穿黑衣，披头散发。这些妇女在我看来很凶猛。但牧者命令她们抬起那些被弃于建筑之外的石头，搬回它们被取来的山上。她们欣然从命，搬走了所有石头，放回它们被取的地方。当所有石头都被移走，塔周围不再有一块石头时，他说：\Quote{我们绕塔走一圈，看看里面是否有缺陷。}于是我和他一起绕塔。牧者见塔建造得美丽，便非常欢喜；因为塔的建造方式，使我见了也羡慕它的建造，它仿佛是由一块石头建造而成，没有一丝接缝。那石头看起来像是从岩石中凿出的，在我看来有如一块独石。\par

% structure: p0089 source=h3 target=subsubsection reason=relative-heading-rank; base=h2

\subsubsection{第十章}

我与他一同行走时，充满了喜悦，目睹如此多美好的事物。牧者对我说：\Quote{你去拿生石灰和精烧粘土来，让我填补那些被搬走又放回建筑的石头所留下的空缺；因为塔的一切都必须平滑。}我便照他的吩咐做了，把东西拿来给他。\Quote{帮我，}他说，\Quote{工作很快就能完成。}于是他填补了那些回到建筑的石头所留下的空缺，命令将塔周围的场地打扫干净；贞女们拿起扫帚打扫，将一切污秽搬出塔外，并提水来，塔周围的地面变得欢快而美丽。牧者对我说：\Quote{一切都已清除；若塔主前来检查，必找不到我们的错。}说完这话，他想要离开；但我抓住他的口袋，开始以主的名义恳求他，请他解释他所显示给我的事。他对我说：\Quote{我必须休息片刻，然后我会向你解释一切；你在这里等我回来。}我对他说：\Quote{先生，我独自一人能做什么？} \Quote{你并非独自一人，}他说，\Quote{因为这些贞女与你同在。} \Quote{那就把我托付给她们吧。}我答道。牧者便召来她们，对她们说：\Quote{我把他托付给你们，直到我回来。}然后离开了。我独自与贞女们在一起；她们很高兴，对我很友善，尤其是其中四位更为卓越的贞女。\par

% structure: p0091 source=h3 target=subsubsection reason=relative-heading-rank; base=h2

\subsubsection{第十一章}

贞女们对我说：\Quote{牧者今天不到这里来。}我说：\Quote{那么，我该做什么？}她们回答说：\Quote{等候他直到他来；如果他要来，他会与你谈话；如果他不来，你就要留在这里和我们一起直到他来。}我对她们说：\Quote{我要等他一直等到天黑；如果他不来，我就进屋去，清早再回来。}她们回答我说：\Quote{你被托付给我们了；你不能离开我们。}我说：\Quote{那么，我该留在哪里？}她们回答说：\Quote{你要像弟兄一样和我们同睡，而不是像丈夫；因为你是我们的弟兄，而且今后我们打算和你住在一起，因为我们极其爱你！}但我羞于和她们一起留下。她们中看来为首的那一个开始吻我。其余的看见她吻我，也开始吻我，并领我绕着塔游玩，与我嬉戏。我也变得像年轻人一样，开始和她们玩耍：有的组成合唱，有的跳舞，有的唱歌；我沉默着，和她们一起绕塔漫步，与她们同乐。天黑时我想进屋去，她们却不让我走，强留了我。于是我和她们一起过夜，睡在塔旁。贞女们把她们的亚麻外衣铺在地上，让我躺在她们中间；她们什么也不做，只是祈祷；我也不停地和她们一起祈祷，并不亚于她们。贞女们因我这样祈祷而欢喜。我和贞女们一直待到第二天第二时辰。然后牧者回来了，对贞女们说：\Quote{你们是否侮辱了他？}她们说：\Quote{问他吧。}我对他说：\Quote{先生，我很高兴和她们在一起。}他问：\Quote{你吃了什么晚饭？}我回答说：\Quote{先生，我整夜享用了主的言语。}他问：\Quote{她们接待你还好吗？}我回答说：\Quote{是的，先生。}他说：\Quote{现在，你希望先听什么？}我说：\Quote{我愿意按你从一开始所示我的次序来听。先生，我求你，我怎样问你，你也怎样给我解释。}他回答说：\Quote{如你所愿，我也要这样给你解释，绝不向你隐瞒任何事。}\par

% structure: p0093 source=h3 target=subsubsection reason=relative-heading-rank; base=h2

\subsubsection{第十二章}

我说：\Quote{首先，先生，请给我解释这一点：这盘石和这门是什么意思？}他回答说：\Quote{这盘石和这门是天主之子。}我说：\Quote{先生，怎么呢？盘石是旧的，门却是新的。}他说：\Quote{听着，无知的人啊，要明白：天主之子比祂所有的受造物更古老，以致祂在创造工程中是圣父的同议者；因此祂是旧的。}我说：\Quote{先生，那么为什么门是新的呢？}他回答说：\Quote{因为祂在救恩计划的末世显现了；因此门被造得新，好让那些要借它得救的人进入天主的国。}他说：\Quote{你看见那些从门进来的石头被用在建造塔上，而那些没有进来的石头又被扔回原处了吗？}我回答说：\Quote{先生，我看见了。}他继续说：\Quote{同样，除非人接受祂的圣名，没有人能进入天主的国。因为如果你想进入一座被城墙环绕、只有一门的城，你能从那城所拥有的门之外进入吗？}我说：\Quote{先生，不然还能怎样呢？}他说：\Quote{如果你不能从城门进入城中，那么同样，人也不能借祂的爱子的名以外的方式进入天主的国。}他接着说：\Quote{你看见那建造塔的群众了吗？}我说：\Quote{先生，我看见了。}他说：\Quote{那些都是光荣的天使，主由此而被围绕。这门是天主之子。这是通往主的唯一入口。因此，人除了通过祂的儿子，无法以其他方式进入祂那里。}他继续说：\Quote{你看见那六个人和中间那位高大光荣的人绕塔行走，并从建筑中抛弃了一些石头吗？}我回答说：\Quote{先生，我看见了。}他说：\Quote{那光荣的人是天主之子，那六位光荣的天使是在祂左右扶持祂的。这些光荣的天使中，没有一位能离开祂而进入天主那里。凡不接受祂的名的人，不能进入天主的国。}\par

% structure: p0095 source=h3 target=subsubsection reason=relative-heading-rank; base=h2

\subsubsection{第十三章}

\Quote{那么，塔}我问，\Quote{是什么意思？} \Quote{这塔}他回答，\Quote{就是教会。} \Quote{这些贞女又是谁？} \Quote{她们是圣神，除非这些贞女将她们的衣袍穿在人身上，人便不能进入天主的国；因为你若只领受那名号，却不从她们领受那衣袍，她们于你毫无裨益。因为这些贞女就是天主子的权能。你若称祂的名号却没有祂的权能，你称祂的名号便是枉然。} \Quote{你所见被弃的那些石头}他继续说，\Quote{曾拥有祂的名号，却没有穿上贞女的衣袍。} \Quote{先生，她们的衣袍是什么样的？}我问。\Quote{那正是她们的名号}他说，\Quote{就是她们的衣袍。凡拥有天主子之名的，也应当拥有这些（贞女）的名号；因为圣子亲自拥有这些贞女的名号。} \Quote{许多石头}他继续说，\Quote{就是你所见的经由这些贞女之手进入塔的建造中并得以存留的，已被穿上她们的力量。因此你看见这塔与那磐石成为一体。同样，凡藉着圣子信了主，并穿上这些圣神的人，将成为一灵、一体，他们衣袍的颜色也将成为一体。而拥有这些贞女名号的人的居所就在这塔内。} \Quote{先生，那些被弃的石头}我问，\Quote{为何被弃？因为它们曾穿过门，并由贞女之手放置于塔的建造中。} \Quote{既然你对一切都有兴趣}他回答，\Quote{并且细致查究，你就听听那些被弃的石头的事。所有这些}他说，\Quote{都领受了天主的名号，也领受了这些贞女的力量。于是他们领受了这些圣神，便变得坚强，与天主的仆人同在；他们拥有一灵、一体、一衣袍。因为他们同心合意，行正义。然而过了一段时间，他们被你所见那些身穿黑衣、肩头裸露、头发散乱且容貌美丽的女子所诱惑。看见这些女子后，他们便想要得到她们，就穿上了她们的力量，脱去了贞女的力量。因此，他们被逐出天主的家，交给了这些女子。而那些没有被这些女子的美貌所迷惑的，仍然留在天主的家。你已得到了被弃之人的解释，}他说。\par

% structure: p0097 source=h3 target=subsubsection reason=relative-heading-rank; base=h2

\subsubsection{第十四章}

\Quote{那么，先生}我说，\Quote{倘若这样的人悔改，除去对那等女子的贪欲，再回到贞女那里，行走于她们的力量与事工中，岂不仍能进入天主的家吗？} \Quote{他们必能进入}他说，\Quote{倘若他们除去那些女子的事工，重新穿上贞女的力量，并行走于她们的事工中。正因如此，建造曾一度停工，好让这些人若悔改，便可进入那塔的建造。但若他们不悔改，就有别人来代替他们，而这些人终将被逐出。为这一切，我感谢了主，因祂怜悯一切呼求祂名号的人；并派遣了补赎天使到我们这些得罪祂的人那里，更新了我们的心灵；当我们已然毁灭，生活再无希望时，祂使我们重获新生。} \Quote{先生，现在我求你}我继续说，\Quote{告诉我为何那塔不是建在平地上，而是建在磐石和门上。} \Quote{你至今仍是}他说，\Quote{无知无悟吗？} \Quote{先生，我必须}我说，\Quote{事事请教您，因为我完全不能明白；这一切伟大而光荣，人类难以理解。} \Quote{请听}他说：\Quote{天主子的名号是伟大的，不可容纳的，且托住整个世界。如果整个受造界都靠天主子托住，那么对于被祂所召、拥有天主子之名、遵守祂诫命的人，你以为如何？你看见祂托住的是怎样的人吗？就是那些以全心拥有祂名号的人。祂自己因此成了他们的根基，并欢欣地托住他们，因为他们不以拥有祂的名号为耻。}\par

% structure: p0099 source=h3 target=subsubsection reason=relative-heading-rank; base=h2

\subsubsection{第十五章}

\Quote{请你，先生，}我说，\Quote{告诉我这些童贞女的名字，以及那些穿着黑色衣服的妇女的名字。} \Quote{听着，}他说，\Quote{那些站在角落里的较强大的童贞女的名字。第一个是信德，第二个是节制，第三个是能力，第四个是忍耐。站在她们中间的其余的童贞女有这些名字：纯朴、无邪、纯洁、喜乐、真理、明悟、和谐、爱德。凡具有这些名字以及天主圣子之名的人，必能进入天国。也听着，}他继续说，\Quote{那些穿着黑色衣服的妇女的名字；其中四个比其他的更强大。第一个是不信，第二个是无节制，第三个是悖逆，第四个是诡诈。她们的随从被称为忧愁、邪恶、淫荡、愤怒、虚谎、愚妄、诽谤、仇恨。天主的仆人若具有这些名字，固然能看见天国，却不能进入其中。} \Quote{先生，那些从坑中取出并被嵌入建筑物的石头，}我说，\Quote{它们是什么？} \Quote{首先，}他说，\Quote{那十块，即被用作地基的，是第一代义人；那二十五块是第二代义人；那三十五块是天主的先知和祂的仆人；那四十块是宗徒和宣讲天主圣子的教师们。} \Quote{那么，先生，}我问，\Quote{为什么童贞女们也把这些石头带过门，并交给她们用于建造塔楼呢？} \Quote{因为，}他回答，\Quote{这些是首先拥有\OriginalTerm{神恩}{spirits}的人，她们彼此从不分离，神恩不离人，人也不离神恩，而是神恩与她们同在，直到她们安眠。如果她们没有这些神恩与她们同在，她们对这座塔楼的建造便毫无用处。}\par

% structure: p0101 source=h3 target=subsubsection reason=relative-heading-rank; base=h2

\subsubsection{第十六章}

\Quote{请再给我解释一点，先生，}我说。\Quote{你想要什么？}他问。\Quote{为什么，先生，}我说，\Quote{这些石头在拥有这些神恩之后，从坑中上来，并被用于塔楼的建造？} \Quote{它们必须，}他回答，\Quote{通过水上来，以便它们能获得生命；因为，除非它们脱下生命的死气，否则无法以任何其他方式进入天国。因此，那些安眠的人也接受了天主圣子的印记。因为，}他继续说，\Quote{在一个人拥有天主圣子之名之前，他是死的；但当他接受印记时，他脱下死气，获得生命。那么，印记就是水：他们死着下到水中，活着上来。因此，这印记也曾向他们宣讲过，他们使用它以便能进入天国。} \Quote{为什么，先生，}我问，\Quote{那四十块石头，虽然已经接受了印记，也同它们一起从坑中上来？} \Quote{因为，}他说，\Quote{这些宗徒和教师宣讲了天主圣子之名，他们在天主圣子的能力和信德中安眠后，不仅向那些安眠的人宣讲，而且自己也给了他们宣讲的印记。因此，他们同他们一起下到水中，又一起上来。【但这些下到水中是活的，又活着上来；而那些先前已安眠的，是死着下去，却活着上来。】因此，借着这些人，他们得以复活并认识天主圣子之名。为此，他们也同他们一起上来，并同他们一起被嵌入塔楼的建造中，未经凿子，就与它们一同被建入。因为他们在正义和极大的纯洁中安眠，只是他们没有这印记。你对此也已经有了解释。}\par

% structure: p0103 source=h3 target=subsubsection reason=relative-heading-rank; base=h2

\subsubsection{第十七章}

\Quote{我明白了，先生，}我回答。\Quote{现在，先生，}我继续说，\Quote{请向我解释关于那些山，为什么它们的形状各不相同、千姿百态。} \Quote{听着，}他说：\Quote{这些山是散居全世界的十二支派。因此，宗徒们曾向它们宣讲了天主圣子。} \Quote{但是为什么这些山有各种不同的形状，有些一个样子，有些另一个样子？请向我解释，先生。} \Quote{听着，}他回答：\Quote{这散居全世界的十二支派是十二个民族。它们在明智和明悟上各不相同。因此，你所看到的山的多样性，也就是这些民族在心思和明悟上的多样性。我将向你解释每一个的行为。} \Quote{首先，先生，}我说，\Quote{请解释这一点：为什么当这些山如此多样时，它们的石头被放入建筑物却变成同一种颜色，像那些从坑中上来的石头一样闪耀发光。} \Quote{因为，}他说，\Quote{普天之下所有的民族，因听见并相信了天主圣子之名而被召唤。因此，接受了印记之后，他们有了同一个明悟和同一个心思；他们的信德成为同一个，他们的爱德也成为同一个，并且随着这名号，他们也带着童贞女的神恩。因此，塔楼的建造变成同一种颜色，明亮如太阳。但是，当他们进入同一地方，成为同一个身体之后，他们中的某些人玷污了自己，被逐出义人的行列，又变成了他们先前的样子，或者更糟。}\par

% structure: p0105 source=h3 target=subsubsection reason=relative-heading-rank; base=h2

\subsubsection{第十八章}

我说道：\Quote{先生，他们认识了天主以后，怎么变得更坏呢？}他回答说：\Quote{那不认识天主而作恶的人，因他的邪恶受到某种惩罚；但那已认识天主的人，不应该再作恶，而应该行善。因此，当他应当行善却作恶时，他岂不是比那不认识天主的人作了更大的恶？因此，那些不认识天主而作恶的人被判处死亡；但那认识天主且见过祂大能作为，却仍继续作恶的人，将受双重惩罚，永远死亡。这样，天主的教会就将被净化。因为你看见那些从塔上被弃、交给恶灵并被逐出的石头，同样，他们也将被逐出，净化后的身体成为一体；正如塔在净化后仿佛成为一块石头。同样，天主的教会也将如此：在它被净化、弃绝了恶人、伪善者、亵渎者、摇摆不定者以及各种作恶者之后，天主的教会将成为一体，一心一意，一智一信一爱。然后天主子将极其喜悦，并为他们欢欣，因为祂接纳了祂的子民是纯洁的。}我说道：\Quote{先生，这一切都是伟大而光荣的。}\par

我说道：\Quote{此外，先生，请向我解释每一座山的力量和行为，使每个灵魂信靠主、听见后，能光荣祂伟大、奇妙而荣耀的名字。}他说：\Quote{请听这众山和十二民族的多样性。}\par

% structure: p0108 source=h3 target=subsubsection reason=relative-heading-rank; base=h2

\subsubsection{第十九章}

\Quote{从第一座山，就是黑色的，那相信的是以下这些人：背教者、亵渎主的人、出卖天主仆人的叛徒。悔改不为他们敞开，死亡却在面前，因此他们也是黑色的，因为他们的种族是无法无天的。从第二座山，就是光秃的，那相信的是以下这些人：伪善者、邪恶的教师。这些人也如同前者，没有公义的果实；因为正如他们的山没有果实，这些人有名无实，信德空虚，在他们内没有真理的果实。他们确有悔改的能力，若他们迅速悔改；但若迟延，他们将与前者一同死亡。}我说道：\Quote{先生，为什么这些人有悔改，而前者却没有？因为他们的行为几乎相同。}他说：\Quote{因此，这些人有悔改，因为他们没有亵渎他们的主，也没有成为天主仆人的叛徒；但因对财富的贪欲，他们成了伪善者，各人按照有罪之人的私欲教导。但他们会受某种惩罚，悔改在他们面前，因为他们不是亵渎者或叛徒。}\par

% structure: p0110 source=h3 target=subsubsection reason=relative-heading-rank; base=h2

\subsubsection{第二十章}

\Quote{从第三座山，就是有荆棘和蒺藜的，那相信的是以下这些人。他们中有一些是富人，另一些深陷于许多事务中。蒺藜是富人，荆棘是那些深陷于许多事务中的人。那些纠缠于许多不同事务的人，不亲近天主的仆人，反而游荡，被他们的交易事务窒息；富人也难以亲近天主的仆人，恐怕这些仆人向他们求什么。因此，这样的人难以进入天主的国。因为正如赤脚走在蒺藜中令人不快，同样，这样的人也难以进入天主的国。但向所有这些人都敞开了悔改，而且是迅速的悔改，以便他们能在这些日子里补足以往未行之事，行些善事，他们便可因天主而生活。但如果他们坚持自己的行为，他们将被交给那些女人，那些女人将处死他们。}\par

% structure: p0112 source=h3 target=subsubsection reason=relative-heading-rank; base=h2

\subsubsection{第二十一章}

\Quote{从第四座山，就是有很多草，上部青绿，根部枯干，有些还被太阳晒焦的，那相信的是以下这些人：怀疑的人，嘴上挂着主，心里却没有主。因此他们的根基枯干，毫无力量；只有他们的话活着，而他们的行为是死的。这样的人是既不活也不死。他们因此就像摇摆不定的人：摇摆不定的人既不枯干也不青绿，既不活也不死。因为正如他们的叶子见到太阳就枯干了，同样，摇摆不定的人听到患难，因恐惧而崇拜偶像，并以他们主的名号为羞。这样的人既不活也不死。但这些人仍可存活，若他们迅速悔改；若不悔改，他们已被交给那些夺去他们生命的女人。}\par

% structure: p0114 source=h3 target=subsubsection reason=relative-heading-rank; base=h2

\subsubsection{第二十二章}

\Quote{从第五座山，就是有青草且崎岖的，那相信的是以下这些人：确实是信徒，但迟钝、固执、自满，想知道一切，却一无所知。因他们的固执，智慧远离了他们，愚蠢无知进入了他们。他们自称有智慧，想成为教师，却缺乏理智。因此因这种傲慢，许多人变得虚空，高抬自己：因为自恃和虚妄的自信是一个大魔鬼。因此，他们中有许多被弃绝，但有些人悔改并相信了，服从那些有智慧的人，认识到自己的愚昧。这类人的其余人仍有悔改的机会；因为他们不是邪恶，而是愚昧无知。因此若这些人悔改，他们将因天主而生活；但若不悔改，他们将在那些行恶的女人中有他们的居所。}\par

% structure: p0116 source=h3 target=subsubsection reason=relative-heading-rank; base=h2

\subsubsection{第二十三章}

\Quote{那些来自第六座山的人——该山有大的和小的裂缝，裂缝中有枯草——信了的人，是以下这些人：占据小裂缝的是那些互相控告的人，因他们的诽谤而在信仰上枯萎了。然而，他们中有许多人悔改了；其余的人当他们听到我的诫命时也将悔改，因为他们的诽谤是小的，他们必迅速悔改。但占据大裂缝的是那些坚持诽谤、在愤怒中彼此报复的人。因此，这些人被拋出塔外，被拒绝参与建塔。这样的人，将难以存活。如果我们的天主和主——他统管万有，且有权柄管理他所有的受造物——不记念那些承认自己罪过的人的邪恶，反而仁慈，那么可朽坏的、充满罪过的人，岂能记念同伴的邪恶，仿佛他能毁灭或拯救他呢？我，这悔改的天使，对你们说：你们中有这种想法的人，当放弃这念头，并悔改，主必医治你们从前的罪，如果你们从这魔鬼中洁净自己；否则，你们必被交给他以致死亡。}\par

% structure: p0118 source=h3 target=subsubsection reason=relative-heading-rank; base=h2

\subsubsection{第二十四章}

\Quote{那些来自第七座山的人——该山草色青翠茂盛，整座山肥沃，各类牲畜和天上的飞鸟都在这山上的草中吃草，而它们所牧放的草越发茂盛——信了的人，是以下这些人：他们总是纯朴、无害、有福，不互相控告，却因天主的仆人而大大喜乐，且被这些贞女的圣神所穿戴，常怜悯众人，用自己的劳动施助众人，毫无责难，毫无迟疑。因此，主看到他们的纯朴和一切的温良，便使他们手所作的工增多，并在他们一切所行的事上赐予恩宠。我，这悔改的天使，对你们这样的人说：愿你们继续如此，你们的后裔必永不被涂抹；因为主曾考验你们，将你们列在我们之中，而你们全体后裔必与天主子同住；因为你们已领受了他的圣神。}\par

% structure: p0120 source=h3 target=subsubsection reason=relative-heading-rank; base=h2

\subsubsection{第二十五章}

\Quote{那些来自第八座山的人——该处有许多泉源，所有天主的受造物都饮于这些泉源——信了的人，是以下这些人：宗徒和教师，他们向普世宣讲，庄严而纯洁地教导主的话语，毫不陷入邪恶的欲望，却常行在正义和真理中，按他们所领受的圣神。因此，这样的人必与众天使一同进入。}\par

% structure: p0122 source=h3 target=subsubsection reason=relative-heading-rank; base=h2

\subsubsection{第二十六章}

\Quote{那些来自第九座山的人——该山荒芜，其中有爬虫和毁灭人的野兽——信了的人，是以下这些人：他们身有污点，作为仆人却尽责不善，掠夺寡妇和孤儿的生活资料，从所接受的职务中为自己牟利。因此，若他们继续受这同样欲望的辖制，他们就是死了，没有生命的希望；但若他们悔改，并圣洁地完成他们的职务，他们便能存活。那些布满疮痂的人，是那些否认了他们的主、未再回到他那里的人；他们变得枯干如旷野，不依附天主的仆人，却独居，毁灭自己的灵魂。正如一棵葡萄树，留在围篱内，被忽略，就毁坏了，被野草荒废，日久变野，对主人不再有用；同样，那些自暴自弃的人，因养成野蛮习惯，对主成为无用。因此，这些人有悔改的能力，除非他们被发现是从心里否认；但若有人被发现是从心里否认，我不知道他是否还能活。我说这话并非为现今的日子，为使任何否认的人能得悔改，因为现在有意否认主的人是不能得救的；但对那些从前否认他的人，悔改似乎是可能的。因此，若有人有意悔改，就当赶快，在塔建成之前；否则，他必被那些妇女彻底毁灭。那些有缺口的石头，是诡诈者和诽谤者；你在第九座山上所见的野兽，正是他们。因为野兽用毒液毁灭杀死一个人，同样，这些人的言语毁灭败坏一个人。因此，他们在信仰上残缺，因他们在自己里面所行的事；然而有些人悔改了，并得了救。其余有这种品性的人，若悔改就能得救；若不悔改，必与那些妇女一同灭亡，他们已承受了她们的能力。}\par

% structure: p0124 source=h3 target=subsubsection reason=relative-heading-rank; base=h2

\subsubsection{第二十七章}

\Quote{那些来自第十座山的人——该处有树木遮蔽某些羊——信了的人，是以下这些人：好客的主教，他们总是毫不虚伪地将天主的仆人迎进家中。这些主教从未停止藉他们的服务保护寡妇和穷乏人，并常保持圣洁的品行。因此，所有这些人都必蒙主永远保护。行这些事的人在天主面前为尊贵，他们的位置已与天使同在，若他们至终事奉天主。}\par

% structure: p0126 source=h3 target=subsubsection reason=relative-heading-rank; base=h2

\subsubsection{第二十八章}

\Quote{从第十一座山，那里有结满果子的树，装饰着各种果实，相信的人是这些：他们为天主子的名字受苦，并且也以全心愉快地受苦，舍弃了自己的生命。} \Quote{为什么，先生，} 我说，\Quote{所有这些树都结果子，有些比其他的更美丽？} \Quote{听着，} 他说：\Quote{所有曾经为主的名受苦的人，在天主面前是可敬的；所有这些人的罪都被赦免了，因为他们为天主子的名字受苦。至于他们的果子为何品种各异，有些更优越，请听。所有，}他继续说，\Quote{那些被带到官府面前受审，没有否认，而是愉快受苦的人——这些人在天主面前更受尊敬，他们的果子也更优越；但所有那些懦弱、犹豫、在心里盘算是否否认或承认，却仍然受苦的人——这些人的果子较少，因为那个念头进入了他们的心；因为那个念头——即仆人应该否认他的主——是邪恶的。因此，你们这些图谋这样事的人要小心，免得那念头留在你们心里，而使你们在天主面前丧亡。你们这些为他的名受苦的人，应当光荣天主，因为他认为你们配得上背负他的名，好使你们的罪都得痊愈。[因此，你们更应自认为有福]，并认为你们做了一件大事，如果你们中有人因天主而受苦。主赐给你们生命，你们却不明白，因为你们的罪孽深重；但如果你们没有为主的名受苦，你们就会因自己的罪而死于天主。我把这些话告诉你们这些犹豫是否否认或承认的人：承认你们有主，免得你们否认他而被交付到监狱。如果外教人责罚他们的奴隶，当其中一个否认他的主人时，你们认为你们的主会做什么？他掌管所有的人。把这些主意从你们心里除掉，好使你们能不断活于天主。}\par

% structure: p0128 source=h3 target=subsubsection reason=relative-heading-rank; base=h2

\subsubsection{第二十九章}

\Quote{从第十二座山（那座白色的山）相信的人是这些：他们如同婴孩，心里没有邪恶产生；也不知道什么是恶，总是像孩子一样。因此，毫无疑问，他们住在天主的国里，因为他们没有玷污天主的任何诫命；反而一生都像孩子一样，保持同样的心志。那么，你们所有将持守不变，并像孩子一样的人，\BibleRef{玛 18:3}，不做邪恶，将比所有先前提到的人更受尊敬；因为所有婴孩在天主面前都是可敬的，且是与他最亲近的人。所以，你们有福了，因为你们从自己身上除去了邪恶，穿上了纯洁。你们将作为最前头的人活于天主。}\par

在他结束了关于山的比喻之后，我对他说：\Quote{先生，现在请向我解释那些从平原取出来、代替从塔中取出的石头而放入建筑中的石头；以及那些被放入建筑的圆石头；还有那些仍然是圆形的石头。}\par

% structure: p0131 source=h3 target=subsubsection reason=relative-heading-rank; base=h2

\subsubsection{第三十章}

\Quote{听着，} 他回答，关于这一切也如此。那些从平原取出来、代替被弃的石头放入塔的建造中的石头，是这座白色山的根基。因此，当所有从白色山相信的人都发现是纯朴时，主命令将从这座山的根基取出的石头投入塔的建造中；因为他知道，如果这些石头进入塔的建造，它们将保持明亮，没有一块变黑。但如果他对其他山也这样决定，他就必须再次巡视那座塔并洁净它。现在，所有相信并且将来会相信的人都被发现是白色的，因为他们属于同一种族。这是有福的种族，因为它是无辜的。现在再听听这些圆而发亮的石头。所有这些也来自白色山。再听听它们为什么是圆的：因为他们的财富使他们从真理上稍微模糊和暗淡了，虽然他们从未离开天主；也没有任何恶言从他们口中发出，而只有一切正义、德行和真理。因此，当主看见这些人的心志，他们生来善良，且能保持善良时，就命令削减他们的财富，不是永远夺走，以便他们能用剩余的做一些好事；他们将活于天主，因为他们属于善良的种族。因此，他们被凿子稍微弄圆，并放入塔的建造中。\par

% structure: p0133 source=h3 target=subsubsection reason=relative-heading-rank; base=h2

\subsubsection{第三十一章}

但那些尚未被用于建造塔楼的圆形石头，也尚未接受印记，它们之所以被放回原处，是因为它们过于圆滑。这个时代必须在这些事情上被削减，并在他们财富的虚空中被削减，然后他们将在天主的国中相遇；因为他们必须进入天主的国，因为主祝福了这个无辜的种族。因此，这个种族中没有人会丧亡；因为即使他们中有人被最邪恶的魔鬼诱惑而犯罪，他也会迅速回到他的主那里。我认为你们有福，我，悔改的使者，你们中凡是像孩童一样无辜的人，因为你们的份是美好的，在天主面前是可敬的。此外，我对你们所有人说，凡接受了天主子印记的，要穿上纯朴，不要记恨，也不要停留在邪恶中。因此，放下对你们冒犯和苦毒的回忆，你们将在一个心神中被塑造。并医治和除去你们中间那些邪恶的分裂，好使群羊之主来时，能因你们而喜乐。祂必喜乐，若祂发现一切健全，你们中没有一人丧亡。但若祂发现这些羊中有一只迷途，牧人有祸了！若牧人自己迷途，他们将为他们的羊群向祂作何答复？他们或许会说，他们被羊群骚扰了？他们不会被相信，因为牧人受羊群之苦是不可信的事；他倒要因他的谎言受罚。而我自己也是个牧人，我负有极其严格的义务，要为你们交账。\par

% structure: p0135 source=h3 target=subsubsection reason=relative-heading-rank; base=h2

\subsubsection{第三十二章}

\Quote{因此，你们要医治自己，趁塔楼还在建造。主住在喜爱和平的人中，因为祂喜爱和平；但祂远离好争竞和极其邪恶的人。因此，要将健全的心神归还给祂，如同你们当初领受的一样。因为当你们将一件新衣服交给洗衣匠，并希望在最后完整地取回它时，如果洗衣匠还给你们一件破衣服，你们会从他那里收下吗？难道不会生气，并责骂他说：\InnerQuote{我给你一件完整的衣服，为什么你把它撕破，使它无用，以致因你所撕的破洞而不能使用？}你们岂不会因在衣服上发现的破洞而对洗衣匠说这一切吗？因此，如果你们为衣服而忧愁，因没有收到完整的衣服而抱怨，你们想主会对你们做什么？祂赐给你们健全的心神，而你们却使它全然无用，以致对拥有者毫无用处。因为它的功用开始变得无益，既然它被你们败坏了。因此，主不会因你们对祂的心神如此行为，而照样对待你们，将你们交付死亡吗？的确，我告诉你们，祂必对一切祂发现仍心怀记恨的人同样对待。不可践踏祂的仁慈，祂说，而要尊敬祂，因为祂对你们的罪如此忍耐，不像你们那样。悔改吧，因为这对你们有益。}\par

% structure: p0137 source=h3 target=subsubsection reason=relative-heading-rank; base=h2

\subsubsection{第三十三章}

\Quote{上面所写的一切，我，牧者，悔改的使者，已向天主的仆人们显示并讲述。因此，如果你们相信并听从我的话，遵行它们，改正你们的道路，你们将能够生活；但如果你们停留在邪恶和记恨中，没有这样的罪人能向天主而生活。我必须说的这些话都已向你们说了。}\par

牧者对说：\Quote{你问了我一切吗？}我回答说：\Quote{是的，先生。}\Quote{你为什么没有问我关于放进建筑的石头形状，好让我解释为什么我们填平了形状？}我说：\Quote{我忘了，先生。}\Quote{那么现在听吧，}他说，\Quote{关于这个也是如此。这些是现在已听见我的诫命，并全心悔改的人。当主看见他们的悔改是良好纯洁的，并且他们能够持守在其中，祂就命令涂抹他们先前的罪。因为这些形状就是他们的罪，它们被削平，好使它们不再显现。}\par

% structure: p0140 source=h2 target=subsection reason=relative-heading-rank; base=h2

\subsection{第十个比喻}

\Emph{论悔改与施舍。}\par

% structure: p0142 source=h3 target=subsubsection reason=relative-heading-rank; base=h2

\subsubsection{第一章}

他于是呼唤我，对我这样说：\Quote{我已把你和你的家交给牧者，使你能被他保护。}\Quote{是的，先生，}我说。\Quote{因此，如果你希望得到保护，}他说，\Quote{免受一切烦恼和一切严酷对待，并在一切善工善言上成功，拥有一切正义的德行，就遵行他给你的这些诫命，你将能制服一切邪恶。因为如果你遵守这些诫命，世界的每一个欲望和欢乐都将服从你，成功将伴随你的一切善工。要接受他的经验和节制，并告诉所有人，他在天主面前享有极大的尊荣和尊严，他是具有大权的掌权者，在他的职位上大有能力。悔改的能力只赐给了他一人，遍及整个世界。他看起来对你有能力吗？但你轻视他的经验，以及他对你的节制。}\par

% structure: p0144 source=h3 target=subsubsection reason=relative-heading-rank; base=h2

\subsubsection{第二章}

我对他说：\Quote{先生，请您亲自问他：自从他进入我家以来，我是否做过任何不当之事，或在任何方面冒犯了他。}他回答说：\Quote{我也知道你既没有做过、也不会做任何不当之事，因此我对你说这些话，是要你坚持到底。因为他曾向我称赞你，你也将把这些话告诉别人，好使那些已经悔改或将要悔改的人能与你怀有同样的心情，而他也就能向我称赞他们，我则向主称赞他们。}我又说：\Quote{先生，我向每一个人传扬天主的伟大作为；并且我盼望所有爱慕这些作为、从前犯过罪的人，听了这些话以后能悔改，重新获得生命。}\Quote{因此，你要继续从事这项职务，并完成它。凡遵行他命令的人，将得生命，并在主那里有极大的尊荣。但那些不遵守他命令的人，便逃避他的生命，藐视他。然而他本人在主那里自有他的尊荣。因此，凡藐视他、不遵行他命令的人，便是自陷于死亡，他们每一个都将为自己的血负责。但我命令你，要服从他的命令，这样你便能医治你从前的罪过。}\par

% structure: p0146 source=h3 target=subsubsection reason=relative-heading-rank; base=h2

\subsubsection{第三章}

\Quote{此外，我派了这些贞女到你这里来，让她们与你同住。因为我看见她们对你和善。你将她们作为助手，好能更好地遵守他的命令：因为没有这些贞女，这些命令是不可能遵守的。我还看见她们乐意留在你家中；但我也要吩咐她们决不离开你的屋子；只要你保持你的屋子洁净，她们就喜欢住在洁净的住所里。因为她们是纯洁、贞洁、勤劳的，在主面前大有影响力。因此，倘若她们发现你的屋子洁净，她们就必与你同在；但若有一点污秽玷污了它，她们就立刻离开你的屋子。因为这些贞女丝毫不喜欢任何污秽。}我对他说：\Quote{先生，我希望我能使她们喜悦，好使她们永远乐意居住在我家中。正如你托付给我的那位并没有抱怨我，她们也不会抱怨我。}他对牧者说：\Quote{我看天主的仆人渴望生活，并遵守这些命令，他必会使这些贞女住进洁净的住所。}他说完这些话，便又把我交托给牧者，并叫了那些贞女来，对她们说：\Quote{既然我看你们乐意住在他家中，我就把他和他的家托付给你们，要求你们决不离开它。}贞女们听了这些话，都很高兴。\par

% structure: p0148 source=h3 target=subsubsection reason=relative-heading-rank; base=h2

\subsubsection{第四章}

天使随后对我说：\Quote{你要在这项职务上刚强行事，向每一个人传扬天主的伟大作为，你必在这项职务上蒙恩。因此，凡行走在这些命令中的人，必得生命，并在他的生活中幸福；但谁若忽视它们，就不得生命，并在今生不幸。你要吩咐所有能够正确行事的人，不要停止行善；因为行善对他们有益。我还要说，每一个人都应当从窘迫中被拯救出来。因为那生活匮乏、每日遭受窘迫的人，处于极大的痛苦和困乏中。因此，凡从这种困乏中拯救这样一个灵魂的人，就为自己赢得了极大的喜乐。因为受这类窘迫折磨的人，所受的痛苦与戴锁链的人相同。此外，许多人因着这类灾祸，因不能忍受，而加速了自己的死亡。因此，谁若知道有这类灾祸折磨一个人而不去救他，就犯了重罪，成为那人之血的罪人。所以，你们这些从主领受了恩惠的人，要行善工；免得因你们迟延不行，致使塔的建造完成，而你们被摒弃于建筑物之外。如今再没有别的塔在建造了。因为建造工程因你们而暂停。除非你们赶紧行善，塔将建成，而你们将被排除在外。}\par

他对我说完这些话后，从卧榻上起身，带着牧者和贞女们离开了。但他对我说，他会把牧者和贞女们再送回我的住所。阿们。\par

\SourceRecord{Translated by F. Crombie.；Ante-Nicene Fathers；Vol. 2.；Edited by Alexander Roberts, James Donaldson, and A. Cleveland Coxe.；Buffalo, NY: Christian Literature Publishing Co.,；1885.；<http://www.newadvent.org/fathers/02013.htm>.}
